\documentclass[letterpaper,11pt,x11names,english,french]{memoir}
  \usepackage{natbib,url,bibentry}
  \usepackage{babel}
  \usepackage[autolanguage]{numprint}
  \usepackage{amsmath,amsthm}
  \usepackage[shortlabels]{enumitem}
  \usepackage{graphicx}
  \usepackage{framed}                  % environnement titled-frame
  \usepackage{manfnt}                  % \mantriangleright (puce)
  \usepackage{dirtree}                 % arbre pour exercice sur \include
  \usepackage{metalogo}                % logo de \XeLaTeX logo
  \usepackage{mflogo}                  % logo de Metafont
  \usepackage{fontawesome}             % plusieurs icônes
  \usepackage{applekeys}               % touches Mac
  \usepackage{answers}                 % exercices et solutions
  \usepackage{listings}                % code informatique
  \usepackage[absolute]{textpos}       % disposition d'images
  \usepackage{pgf}                     % transparence pour couverture avant
  \usepackage{ifthen}                  % exécution conditionnelle


  %%% =============================
  %%%  Informations de publication
  %%% =============================
  \renewcommand{\year}{2016}
  \renewcommand{\month}{11-3}
  \newcommand{\ISBN}{978-2-9811416-7-5}
  \newcommand{\ctanurl}{https://ctan.org/pkg/formation-latex-ul/}

  %%% ===================
  %%%  Style du document
  %%% ===================

  %% Polices de caractères
  \usepackage{fontspec}
  \usepackage[bold-style=ISO]{unicode-math}
  \defaultfontfeatures{Scale=0.92}
  \setmainfont[Ligatures=TeX,Numbers=OldStyle]{Lucida Bright OT}
  \setmathfont{Lucida Bright Math OT}
  \setmonofont{Lucida Grande Mono DK}
  \setsansfont[Scale=1.0,Numbers=OldStyle]{Myriad Pro}
  \newfontfamily\fullcaps[Letters=Uppercase,Numbers=Uppercase]{Myriad Pro}
  \usepackage[babel=true]{microtype}
  \usepackage{icomma}

  %% Polices additionnelles pour le chapitre trucs et astuces
  \newfontfamily\CM{cmunrm.otf}                       % Computer Modern
  \newfontfamily\Times{texgyretermes-regular.otf}     % Times
  \newfontfamily\Palatino{texgyrepagella-regular.otf} % Palatino
  \newfontfamily\Bookman{texgyrebonum-regular.otf}    % Bookman
  \newfontfamily\NewCent{texgyreschola-regular.otf}   % New Cent. Sch.
  \newfontfamily\Charter{XCharter-Roman.otf}          % Charter
  \newfontfamily\Helvet{texgyreheros-regular.otf}     % Helvetica

  %% Couleurs
  \usepackage{xcolor}
  \definecolor{comments}{rgb}{0.7,0,0}    % rouge foncé
  \definecolor{link}{rgb}{0,0.4,0.6}      % ~RoyalBlue de dvips
  \definecolor{url}{rgb}{0.6,0,0}         % rouge-brun
  \definecolor{citation}{rgb}{0,0.5,0}    % vert foncé
  \definecolor{ULlinkcolor}{rgb}{0,0,0.3} % de ulthese.cls
  \definecolor{rouge}{rgb}{0.85,0,0.07}   % rouge bandeau identitaire
  \definecolor{or}{rgb}{1,0.8,0}          % or bandeau identitaire

  %% Hyperliens
  \usepackage{hyperref}
  \hypersetup{%
    pdfauthor = {Vincent Goulet},
    pdftitle = {Rédaction avec LaTeX},
    colorlinks = {true},
    linktocpage = {true},
    urlcolor = {url},
    linkcolor = {link},
    citecolor = {citation},
    pdfpagemode = {UseOutlines},
    pdfstartview = {Fit},
    bookmarksopen = {true},
    bookmarksnumbered = {true},
    bookmarksdepth = {subsubsection}}
  \setlength{\XeTeXLinkMargin}{1pt}

  %% Étiquettes de \autoref (redéfinitions compatibles avec babel).
  %% Attention! Les % à la fin des lignes sont importants sinon des
  %% blancs apparaissent dès que la commande \selectlanguage est
  %% utilisée, notamment dans la bibliographie.
  \addto\extrasfrench{%
    \def\subsectionautorefname{section}%
    \def\figureautorefname{figure}%
    \def\tableautorefname{tableau}%
    \def\exempleautorefname{exemple}%
    \def\exerciceautorefname{exercice}%
    \def\appendixautorefname{annexe}%
  }

  %% Table des matières (inspirée de classicthesis.sty)
  \renewcommand{\cftchapterleader}{\hspace{1.5em}}
  \renewcommand{\cftchapterafterpnum}{\cftparfillskip}
  \renewcommand{\cftsectionleader}{\hspace{1.5em}}
  \renewcommand{\cftsectionafterpnum}{\cftparfillskip}

  %% Titres des chapitres
  \chapterstyle{hangnum}
  \renewcommand{\chaptitlefont}{\normalfont\Huge\sffamily\bfseries\raggedright}

  %% Marges, entêtes et pieds de page
  \setlength{\marginparsep}{7mm}
  \setlength{\marginparwidth}{20mm}
  \setlength{\headwidth}{\textwidth}
  \addtolength{\headwidth}{\marginparsep}
  \addtolength{\headwidth}{\marginparwidth}
  \addtolength{\marginparwidth}{15mm} % plus d'espace pour titres de documentation

  %% Titres des sections et sous-sections
  \setsecheadstyle{\normalfont\Large\sffamily\bfseries\raggedright}
  \setsubsecheadstyle{\normalfont\large\sffamily\bfseries\raggedright}
  \maxsecnumdepth{subsection}
  \setsecnumdepth{subsection}

  %% Listes. Paramétrage avec enumitem.
  \setlist[enumerate]{leftmargin=*,align=left}
  \setlist[enumerate,2]{label=\alph*),labelsep=*,leftmargin=1.5em}
  \setlist[enumerate,3]{label=\roman*),labelsep=*,leftmargin=1.5em,align=right}
  \setlist[itemize]{leftmargin=*,align=left}

  %% Paramétrage de babel
  \frenchbsetup{%
    StandardItemizeEnv=true,       % format standard des listes
    ThinSpaceInFrenchNumbers=true, % espace fine dans les nombres
    ItemLabeli=\mantriangleright,  % puce premier niveau
    ItemLabelii=\textendash,       % puce second niveau
    og=«, fg=»                     % caractères « et » sont les guillemets
  }
  \def\frenchfigurename{{\scshape Fig.}}
  \def\frenchtablename{{\scshape Tab.}}

  %% Sections de code source
  \lstloadlanguages{[LaTeX]TeX}
  \lstset{language=[LaTeX]TeX,
    basicstyle=\ttfamily\NoAutoSpacing,
    keywordstyle=\mdseries,
    commentstyle=\color{comments}\slshape,
    emphstyle=\bfseries,
    escapeinside=`',
    extendedchars=true,
    showstringspaces=false,
    backgroundcolor=\color{LightYellow1},
    frame=lr, rulecolor=\color{LightYellow1},
    xleftmargin=3.4pt, xrightmargin=3.4pt}

  %%% =========================
  %%%  Nouveaux environnements
  %%% =========================

  %% Exemples
  \theoremstyle{definition}
  \newtheorem{exemple}{Exemple}[chapter]

  %% Exercices et réponses
  \Newassociation{sol}{solution}{solutions}
  \newcounter{exercice}[chapter]
  \renewcommand{\theexercice}{\thechapter.\arabic{exercice}}
  \newenvironment{exercice}[1][]{%
    \begin{list}{}{%
        \refstepcounter{exercice}
        \ifthenelse{\equal{#1}{nosol}}{%
          \renewcommand{\makelabel}{\bfseries\theexercice}}{%
          \hypertarget{ex:\theexercice}{}
          \Writetofile{solutions}{\protect\hypertarget{sol:\theexercice}{}}
          \renewcommand{\makelabel}{%
            \bfseries\protect\hyperlink{sol:\theexercice}{\theexercice}}}
        \settowidth{\labelwidth}{\bfseries\theexercice}
        \setlength{\leftmargin}{\labelwidth}
        \addtolength{\leftmargin}{\labelsep}
        \setlist[enumerate,1]{label=\alph*),labelsep=*,leftmargin=1.5em}
        \setlist[enumerate,2]{label=\roman*),labelsep=0.5em,align=right}}
      \item}
    {\end{list}}
  \renewenvironment{solution}[1]{%
    \begin{list}{}{%
        \renewcommand{\makelabel}{%
          \bfseries\protect\hyperlink{ex:#1}{#1}}
        \settowidth{\labelwidth}{\bfseries #1}
        \setlength{\leftmargin}{\labelwidth}
        \addtolength{\leftmargin}{\labelsep}
        \setlist[enumerate,1]{label=\alph*),labelsep=*,leftmargin=1.5em}
        \setlist[enumerate,2]{label=\roman*),labelsep=0.5em,align=right}}
      \item}
    {\end{list}}

  %% Démo de code LaTeX. Le code de 'texample' et 'eqxample' est
  %% repris de amsldoc.tex avec des petits ajustements.
  \newenvironment{demo}{%
    \begin{trivlist}\item}{%
    \end{trivlist}}
  \newenvironment{texample}[1][0.5\linewidth]{%
    \noindent\begin{minipage}{#1}%
      \def\producing{\end{minipage}\hfill\begin{minipage}{\dimexpr0.97\linewidth-#1}%
        \hbox\bgroup\kern-.2pt%
        \vbox\bgroup\parindent0pt\relax
        % The 3pt is to cancel the -\lineskip from \displ@y
        \abovedisplayskip3pt \abovedisplayshortskip\abovedisplayskip
        \belowdisplayskip0pt \belowdisplayshortskip\belowdisplayskip
        \noindent}
    }{%
      \par
      % Ensure that a lonely \[\] structure doesn't take up width less than
      % \hsize.
      \hrule height0pt width\hsize
      \egroup\kern-.2pt\egroup
    \end{minipage}%
    \par
  }
  \newenvironment{eqxample}{%
    \noindent\begin{minipage}{.5\columnwidth}%
      \def\producing{\end{minipage}\hfill\begin{minipage}{.45\columnwidth}%
        \hbox\bgroup\kern-.2pt\vrule width.2pt%
        \vbox\bgroup\parindent0pt\relax
        % The 3pt is to cancel the -\lineskip from \displ@y
        \abovedisplayskip3pt \abovedisplayshortskip\abovedisplayskip
        \belowdisplayskip0pt \belowdisplayshortskip\belowdisplayskip
        \noindent}
    }{%
      \par
      % Ensure that a lonely \[\] structure doesn't take up width less than
      % \hsize.
      \hrule height0pt width\hsize
      \egroup\vrule width.2pt\kern-.2pt\egroup
    \end{minipage}%
    \par
  }

  %% Un exemple du chapitre Trucs et astuces nécessite des
  %% environnements 'lstlisting' imbriqués, ce que ne digère pas
  %% LaTeX. La ruse consiste à définir un environnement équivalent qui
  %% porte simplement un autre nom.
  \lstnewenvironment{vglisting}{\lstset{deletetexcs={int,include}}}{}

  %% Exemples de notices bibliographiques
  \newenvironment{bibexample}[1][\linewidth]{%
    \begin{minipage}[t]{#1}%
      \begin{trivlist}}
      {\end{trivlist}\end{minipage}}

  %% Encadré générique pour les remarques importantes et autres
  %% comportant une icône sur la gauche. Argument: symbole à
  %% placer sur la gauche (obligatoire).
  \newenvironment{infobox}[1]{%
    \setlength{\FrameRule}{1pt}
    \begin{table}[h]%
      \begin{framed}%
        \noindent
        \begin{minipage}{0.1\linewidth}
          \raisebox{-1.5em}[0em][0em]{\HUGE#1}
        \end{minipage}
        \begin{minipage}[t]{0.88\linewidth}}%
        {\end{minipage}\end{framed}\end{table}}

  %% Remarques importantes
  \newenvironment{important}{%
    \begin{infobox}{\faExclamationCircle}}%
    {\end{infobox}}

  %% Informations
  \newenvironment{information}{%
    \begin{infobox}{\faInfoCircle}}%
    {\end{infobox}}

  %% Encadré avec titre (basé sur 'titled-frame' de framed pour les
  %% conseils du TeXpert. Cet environnement est laissé flottant.
  \newenvironment{conseil}{%
    \colorlet{TFFrameColor}{black}%
    \colorlet{TFTitleColor}{white}%
    \begin{table}%
      \begin{titled-frame}{\sffamily Conseil du {\TeX}pert}%
        \noindent
        \begin{minipage}{0.1\linewidth}
          \raisebox{-1.5em}[0em][0em]{\HUGE\faThumbsOUp}
        \end{minipage}
        \begin{minipage}[t]{0.88\linewidth}}%
        {\end{minipage}\end{titled-frame}\end{table}}

  %%% =======
  %%%  Index
  %%% =======
  \renewcommand{\preindexhook}{%
    Cet index contient des références aux commandes et environnements
    {\LaTeX}, ainsi qu'aux noms de paquetages et de classes. %
    Le premier numéro indique habituellement, mais pas toujours,
    la page où un concept est introduit, défini ou expliqué.%
    \vskip\onelineskip}
  \lstset{language=[AlLaTeX]TeX,
    morekeywords={align,align*,aligned,bmatrix,cases,equation*,%
      figure,gather,lstlisting,multline,quote,split,%
      table,tabular,tabularx},
    deletekeywords={document},   % répéter dans deletetexcs
    moretexcs={toprule,midrule,bottomrule,%
      includegraphics,reflectbox,resizebox,rotatebox,scalebox,%
      includepdf,frenchfigurename,frenchtablename,%
      newsubfloat,subcaption,%
      bm,dfrac,tfrac,iint,text,mathcal,mathbb,eqref,symbf,%
      citet,citep,citeauthor,citeyear,%
      setmainfont,setsansfont,setmonofont,setmathfont,%
      color,textcolor,definecolor,colorlet,hypersetup},
    deletetexcs={document,documentclass,usepackage,begin,end,LaTeX,TeX,%
      normalfont,bfseries,textbf,itshape,scshape,sffamily,ttfamily,texttt,%
      emph,small,Huge,raggedright,%
      hfill,def,a,b,c,d,em,i,j,l,r,t},
    index=[1][keywords],        % environnements
    indexstyle=[1]\ixenv,
    index=[2][texcs],           % commandes
    indexstyle=[2]\ixcmd}
  \newcommand{\ixenv}[1]{\index{#1 env@\Pe{#1} (environnement)}%
    \index{environnement!#1@\Pe{#1}}}
  \newcommand{\ixcmd}[1]{\index{#1@\string\cs{#1}}}
  \makeindex

  %%% =====================
  %%%  Nouvelles commandes
  %%% =====================

  %% Noms de fonctions, code, environnement, etc.
  \newcommand{\code}[1]{\texttt{#1}}
  \newcommand{\fichier}[1]{\texttt{#1}}
  \newcommand{\class}[1]{\textsf{#1}\index{#1 class@\textsf{#1} (classe)}%
    \index{classe!#1}}
  \newcommand{\pkg}[1]{\textbf{#1}\index{#1 pkg@\textbf{#1} (paquetage)}%
    \index{paquetage!#1}}
  \newcommand{\Pe}[1]{\code{#1}}         % tiré de la doc de memoir
  \newcommand{\Ie}[1]{\Pe{#1}\ixenv{#1}} % idem
  \newcommand{\mat}[1]{\symbf{#1}}       % en mode mathématique

  %% Modification de commandes tirées de memoir.cls servant à afficher
  %% des noms de commandes.
  %% - \cmdprint est modifiée pour que le nom de la commande ne soit
  %%   pas en italique;
  %% - \cmd est modifiée pour utiliser @ comme séparateur dans \index
  %%   et pour utiliser \cs plutôt que \cmdprint pour afficher le nom de
  %%   la commande (afin d'obtenir le même format d'entrée d'index
  %%   qu'avec \ixcmd ci-dessus);
  %% - \pixbsbs et \pixabang servent respectivement à afficher et
  %%   indexer \\ et \! ;
  %% - \pixbar sert à afficher et indexer \| avec un hack pour
  %%   contourner un problème d'insertion de l'hyperlien vers le
  %%   numéro de page: l'entrée est triée sur le symbole [ plutôt que
  %%   sur |.
  \renewcommand{\cmdprint}[1]{\textup{\texttt{\string#1}}}
  \makeatletter
  \renewcommand{\cmd}[1]{\cmdprint{#1}%
    \index{\expandafter\@gobble\string#1@\string\cs{\expandafter\@gobble\string#1}}}
  \makeatother
  \newcommand*{\pixbsbs}{\cmdprint{\\}\index{"\ @\string\cs{}\bs}}
  \newcommand*{\pixabang}{\cmdprint{\!}\index{"!@\string\cs{}\texttt{"!}}}
  \newcommand*{\pixbar}{\cmdprint{\|}\index{[@\string\cs{}\texttt{\textbar}}}

  %% Indications de capsule vidéo
  \newcommand{\capsule}[2]{\href{#1}{#2~\raisebox{-0.2ex}{\Large\faYoutubePlay}}}

  %% Hyperlien avec symbole de lien externe juste après
  \newcommand{\link}[2]{\href{#1}{#2~\raisebox{-0.2ex}{\faExternalLink}}}

  %% Lien vers documentation dans la marge
  %% usage: \doc[documentation]{nom_fichier}{url}
  \newcommand{\doc}[3][documentation]{\link{#3}{#1}%
    \ifthenelse{\equal{#2}{}}{}{\marginpar%
      [\hfill\faBookmark~\fichier{#2}]%
      {\faBookmark~\fichier{#2}}}}

  %% Suppression de l'hyperlien
  \newcommand{\nolink}[1]{\begin{NoHyper}#1\end{NoHyper}}

  %% «Bouton» de la page de copyright
  \newcommand{\browsebutton}{%
    \setlength{\fboxrule}{1pt}%
    \framebox[40mm][c]{%
      \rule[-5pt]{0mm}{16pt}%
      \makebox[7mm]{\raisebox{-2pt}{\LARGE\faGithub}}\;%
      {\sffamily Voir sur GitHub}}}

  %% Pour le tableau des commandes d'espacement en mode mathématique.
  %% Pris de la doc de amsmath.
  \newcommand{\lspx}{\mathord{\dashv\mkern-3mu}}
  \newcommand{\rspx}{\mathord{\mkern-2mu\vdash}}
  \newcommand{\spx}[1]{$\lspx #1\rspx$}

  %% Logo BIBTeX; tiré de http://bit.ly/1RQqUnG
  \newcommand{\BibTeX}{\rmfamily B\kern-.05em{\scshape i\kern-.025em b}\kern-.08em \TeX}

  %% Chapitre sur les bibliographies: des références bibliographiques
  %% sont insérées dans le texte avec \bibentry. Certaines commandes
  %% de francaisbst.tex sont alors utilisées, mais non encore
  %% définies. Répétées ici. De plus, il faut définir ici la commande
  %% \enquote plutôt que dans francais.bst. C'est pourquoi il y a une
  %% version modifiée de ce fichier dans ces sources.
  %% Voir http://bit.ly/1MORZmp
  \global\def\bbland{et}
  \global\def\bbledn{\'ed.}
  \global\def\bblfourtho{4{\ieme}}
  \global\def\bblth{{\ieme}}
  \global\def\bblvol{vol.}
  \def\bblno{\no{}}
  \def\bblpp{p.}
  \newcommand{\enquote}[1]{\guillemotleft#1\guillemotright}

  %%% =======
  %%%  Varia
  %%% =======

  %% Style de la bibliographie
  \bibliographystyle{francais}

  %% Longueurs pour la composition des pages couvertures avant et
  %% arrière.
  \newlength{\banderougewidth} \newlength{\banderougeheight}
  \newlength{\bandeorwidth}    \newlength{\bandeorheight}
  \newlength{\imageheight}
  \newlength{\logoheight}
  \newlength{\gapwidth}


  %%% ============
  %%%  Page titre
  %%% ============
  \title{\protect\raggedright%
    \sffamily\bfseries
    \fontsize{40}{40}\selectfont
    Rédaction avec \\
    \rmfamily\mdseries
    \fontsize{45}{45}\selectfont
    \raisebox{10pt}{{\textbackslash}title}%
    \fontsize{80}{80}\selectfont%
    \{%
    \fontsize{70}{70}\selectfont%
    \LaTeX
    \fontsize{80}{80}\selectfont%
    \}}
  \author{\protect\raggedright%
    \sffamily\bfseries
    \fontsize{17}{20}\selectfont
    Vincent Goulet \\
    \mdseries
    \fontsize{15}{20}\selectfont
    Professeur titulaire \textbar\ École d'actuariat}
  \date{%
    \sffamily\mdseries\fontsize{15}{20}\selectfont
    Édition {\fullcaps\year}.\month}

%  \includeonly{couverture-avant,frontispice}

\begin{document}

\frontmatter

%% Page couverture avant.
\pagestyle{empty}
\begingroup

\TPGrid{3}{36}
\textblockorigin{0mm}{0mm}
\setlength{\parindent}{0mm}
\setlength{\banderougewidth}{2\TPHorizModule}
\setlength{\banderougeheight}{\TPVertModule}
\setlength{\bandeorwidth}{\TPHorizModule}
\setlength{\bandeorheight}{\banderougeheight}
\setlength{\imageheight}{29\TPVertModule}
\setlength{\logoheight}{2.5\TPVertModule}
\setlength{\gapwidth}{1.5pt}
\addtolength{\bandeorwidth}{-\gapwidth}
\addtolength{\imageheight}{-\gapwidth}

%% bandeau identitaire
\begin{textblock*}{8.5in}[0,1](0mm,30\TPVertModule)
  \textcolor{rouge}{\rule{\banderougewidth}{\banderougeheight}}% % bande rouge
  \rule{\gapwidth}{0pt}%                                         % filet
  \textcolor{or}{\rule{\bandeorwidth}{\bandeorheight}}           % bande or
\end{textblock*}

%% logo UL
\begin{textblock*}{\TPHorizModule}(2\TPHorizModule,31\TPVertModule)
  \rule{\gapwidth}{0pt}%                                     % filet
  \includegraphics[height=\logoheight,keepaspectratio=true]{ul_p}
\end{textblock*}

%% image de fond
\begin{textblock*}{8.5in}(0mm,0mm)
  \includegraphics[height=\imageheight,keepaspectratio=true]{Suricata.jpg}
\end{textblock*}

%% trame
\begin{textblock*}{2\TPHorizModule}(0mm,13\TPVertModule)
  \pgfsetfillopacity{0.5}
  \textcolor{white}{\rule{\linewidth}{8\TPVertModule}}
  \pgfsetfillopacity{1}
\end{textblock*}

%% titre
\begin{textblock*}{1.7\TPHorizModule}(0.3\TPHorizModule,14.2\TPVertModule)
  \thetitle
\end{textblock*}

\endgroup

%%% Local Variables:
%%% mode: latex
%%% TeX-master: "formation-latex-ul"
%%% TeX-engine: xetex
%%% End:

\null\cleardoublepage           % cf. section 2.2 textpos.pdf

%% Page frontispice
\begingroup

\TPGrid{3}{36}
\textblockorigin{0mm}{0mm}

%% auteur
\begin{textblock*}{1.7\TPHorizModule}(0.3\TPHorizModule,5\TPVertModule)
  \begin{minipage}{\textwidth}
    \theauthor
  \end{minipage}
\end{textblock*}

%% titre
\begin{textblock*}{1.7\TPHorizModule}(0.3\TPHorizModule,14.2\TPVertModule)
  \begin{minipage}{\textwidth}
    \thetitle \\
  \end{minipage}
\end{textblock*}

%% édition
\begin{textblock*}{1.7\TPHorizModule}(0.3\TPHorizModule,30\TPVertModule)
  \begin{minipage}{\textwidth}
    \thedate
  \end{minipage}
\end{textblock*}

\endgroup

%%% Local Variables:
%%% mode: latex
%%% TeX-master: "formation-latex-ul"
%%% TeX-engine: xetex
%%% End:

\null\clearpage                 % idem

%% Page de copyright
version https://git-lfs.github.com/spec/v1
oid sha256:5cefc45a96685bc8ee9467d1c8680da0b2962093c874b037bc0cf76c3357be6b
size 2648

\clearpage

%% Corps du document
\pagestyle{companion}

version https://git-lfs.github.com/spec/v1
oid sha256:9aa4ffc74913f8ccf218e3e5d4fc7f18d6c6f9899c18ef5bc6ee16648a338c94
size 5162


\cleartorecto
\tableofcontents*

\mainmatter

version https://git-lfs.github.com/spec/v1
oid sha256:4e3246cb2f0de69af95c53244a7bc76c6cfc69807e82331022f7017b660f8153
size 13850

\chapter{Principes de base}
\label{chap:bases}

Un système de mise en page tel que {\LaTeX} repose sur une logique de
séparation entre l'apparence d'un document et sa structure. La
personne habituée à utiliser un traitement de texte devra fort
probablement se défaire d'une vilaine habitude: se préoccuper sans
cesse, au moment de la rédaction, de la disposition du texte.

Ce principe accepté, il faudra néanmoins indiquer au logiciel la
structure du document. Avec {\LaTeX} cela s'effectue par le biais de
diverses instructions que l'on insère au fil du texte. À la base, les
logiciels de traitements de texte n'opèrent pas différemment, sauf
qu'ils cachent les codes aux utilisateurs\footnote{%
  Le leader du traitement de texte jusqu'au milieu des années 1990,
  Word~Perfect, offrait l'option d'afficher l'ensemble des codes de
  mise en page. C'est malheureusement une caractéristique brillante
  que Microsoft Word et les autres progiciels développés depuis ont
  choisi d'omettre.}. %
Sous prétexte de simplicité d'utilisation, ils causent en fait bien
des mots de tête. Qui ne s'est pas déjà demandé, en utilisant un
traitement de texte moderne: «Pourquoi donc mon texte est-il
soudainement en gras?»

Ce chapitre explique comment aborder la rédaction d'un document avec
{\LaTeX} ainsi que la syntaxe de base des différents types
d'instructions que l'on peut insérer dans un texte pour en spécifier
la structure et la mise en forme.


\section{Séparation du contenu et de l'apparence}
\label{sec:bases:separation}

Lors de la rédaction avec un système de mise en page tel que {\LaTeX},
on se concentre sur le contenu et la \emph{structure} du document, et
non pas sur son \emph{apparence}. Par exemple:
\begin{itemize}
\item au lieu de prescrire qu'un titre de section doit être en gras
  14~points, on indique simplement à {\LaTeX} que le texte doit être
  traité comme un titre de section;
  \begin{demo}
    \begin{minipage}{0.45\linewidth}
\begin{lstlisting}
\textbf{\large Titre}
\end{lstlisting}
    \end{minipage}
    \hfill \faArrowRight \hfill
    \begin{minipage}{0.45\linewidth}
\begin{lstlisting}
\section{Titre}
\end{lstlisting}
    \end{minipage}
  \end{demo}
\item au lieu de décider qu'un mot sur lequel l'on souhaite insister
  sera en italique, on indique à {\LaTeX} de mettre de l'emphase sur
  ce mot sans se soucier de la mise en forme.
  \begin{demo}
    \begin{minipage}{0.45\linewidth}
\begin{lstlisting}
\textit{texte}
\end{lstlisting}
    \end{minipage}
    \hfill \faArrowRight \hfill
    \begin{minipage}{0.45\linewidth}
\begin{lstlisting}
\emph{texte}
\end{lstlisting}
    \end{minipage}
  \end{demo}
\end{itemize}

L'apparence du texte sera prise en charge par {\LaTeX}. Comme les
gabarits sont l'{\oe}uvre de spécialistes en typographie, il est
généralement préférable de ne pas les modifier. À titre d'exemple,
{\LaTeX} détermine automatiquement la largeur des marges en fonction
de la taille de la police de manière à ce que les lignes de texte
comptent approximativement 70~caractères. La raison: lorsqu'une ligne
de texte est trop longue, notre {\oe}il a de la difficulté à la suivre
sur toute sa longueur. Il a tendance à passer à la ligne inférieure,
ce qui rend la lecture plus difficile.


\section{Règles de saisie}
\label{sec:bases:saisie}

Une fois le principe de séparation du contenu et de l'apparence
compris et accepté, on veillera, lors de la saisie du texte, à
respecter les règles simples suivantes.
\begin{enumerate}
\item On sépare les mots par une ou plusieurs \emph{espaces}. Qu'il y
  en ait une ou un millier, seule la première compte et la mise en
  page sera la même.
  \begin{demo}
    \begin{texample}
\begin{lstlisting}
Les espaces délimitent les
mots. Leur nombre n'a pas
d'importance.
\end{lstlisting}
      \producing
      Les espaces délimitent les
      mots. Leur nombre n'a pas
      d'importance.
    \end{texample}
    \begin{texample}
\begin{lstlisting}[showstringspaces=true]
Les  espaces   délimitent
les                 mots.
Leur    nombre  n'a pas
d'importance.
\end{lstlisting}
      \producing
      Les  espaces   délimitent
      les                 mots.
      Leur    nombre  n'a pas
      d'importance.
    \end{texample}
  \end{demo}
%
\item On sépare les paragraphes par une ou plusieurs lignes blanches.
  Celles-ci n'apparaîtront pas nécessairement dans le texte final; les
  gabarits standards identifient les paragraphes par un retrait de
  première ligne.
  \begin{demo}
    \begin{texample}
\begin{lstlisting}
Les lignes blanches
délimitent les paragraphes.

Une ou plusieurs, ça ne fait
aucune différence!
\end{lstlisting}
      \producing
Les lignes blanches
délimitent les
paragraphes.

        Une ou plusieurs, ça ne fait aucune différence!
    \end{texample}
\begin{texample}
\begin{lstlisting}
Les lignes blanches
délimitent les paragraphes.



Une ou plusieurs, ça ne fait
aucune différence!
\end{lstlisting}
      \producing
        Les lignes blanches  délimitent
        les paragraphes.



        Une ou plusieurs, ça ne fait aucune différence!
    \end{texample}
  \end{demo}
%
\item On utilise des \emph{commandes} pour indiquer la structure du
  texte dans le code source. Celles-ci débutent presque toujours par
  le symbole {\bs}. À la différence des logiciels de traitement de
  texte, les instructions de mise en forme du document sont donc
  toujours visibles et, par conséquent, modifiables facilement et sans
  surprise (on ne se demande donc jamais où termine le gras).
  \begin{demo}
    \begin{texample}
\begin{lstlisting}
Les commandes sont visibles
dans le \textbf{code source},
mais évidemment pas dans le
\emph{document} fini.
\end{lstlisting}
      \producing
      Les commandes sont visibles
      dans le \textbf{code source},
      mais évidemment pas dans le
      \emph{document} fini.
    \end{texample}
  \end{demo}
\end{enumerate}


\section{Structure d'un fichier}
\label{sec:bases:structure}

Un fichier source {\LaTeX} --- dont on trouvera un exemple simple à la
\autoref{fig:bases:parties} --- est toujours composé de deux parties.

\begin{figure}
  \centering
  \begin{minipage}{0.75\linewidth}
\begin{lstlisting}[numbers=left, numberstyle=\tiny,
                   frame=single, rulecolor=\color{black}, framesep=6pt]
\documentclass[11pt,french]{article} `\label{lst:bases:preambule_debut}'
  \usepackage{babel}
  \usepackage[autolanguage]{numprint}
  \usepackage[utf8]{inputenc}
  \usepackage[T1]{fontenc} `\label{lst:bases:preambule_fin}'

\begin{document} `\label{lst:bases:corps_debut}'

Lorem ipsum dolor sit amet, consectetur
adipiscing elit. Donec quam nulla, bibendum
vitae ipsum vel, fermentum pellentesque orci.

\end{document} `\label{lst:bases:corps_fin}'
\end{lstlisting}
  \end{minipage}
  \caption{Fichier source {\LaTeX} simple comportant les deux parties
    obligatoires: le préambule (lignes
    \ref*{lst:bases:preambule_debut}--\ref*{lst:bases:preambule_fin})
    et le corps du document (lignes
    \ref*{lst:bases:corps_debut}--\ref*{lst:bases:corps_fin}).}
  \label{fig:bases:parties}
\end{figure}

\begin{description}
\item[Préambule] Suite de commandes spécifiant la mise en forme
  globale du document (format du papier, marges, entête et pied de
  page, etc.). Il contient au minimum la commande
  \cmd{\documentclass}. Les commandes contenues dans le préambule ont
  un effet global sur le document.

  Dans l'exemple de la \autoref{fig:bases:parties}, le préambule
  s'étend de la ligne \ref*{lst:bases:preambule_debut} à la ligne
  \ref*{lst:bases:preambule_fin}.
\item[Corps du document] Contenu du document en tant que tel.
  Il débute par \verb=\begin{document}= et se termine par
    \verb=\end{document}=. Le corps du document peut aussi contenir
  des commandes, mais l'effet de celles-ci demeure presque toujours local.

  Les lignes \ref*{lst:bases:corps_debut}--\ref*{lst:bases:corps_fin}
  du code de la \autoref{fig:bases:parties} forment le corps du
  document.
\end{description}


\section{Classes et paquetages}
\label{sec:bases:classes}

La première commande du préambule est normalement la déclaration de la
\emph{classe} du document. La forme de la déclaration est la suivante:
\begin{lstlisting}
\documentclass`\oarg{options}\marg{classe}'
\end{lstlisting}
Les classes standards de {\LaTeX} sont \class{article},
\class{report}, \class{book}, \class{letter} et \class{slides}. La
classe pour les thèses et mémoires de l'Université Laval se nomme
\class{ulthese} \citep{ulthese}. Celle-ci se base sur la classe
\class{memoir} et, par conséquent, elle hérite de toutes ses (nombreuses)
fonctionnalités; consulter l'\autoref{chap:ulthese}. La
\autoref{sec:organisation:classe} traite des différences entre les
diverses classes et des \meta{options} disponibles.

Les \emph{paquetages} permettent de modifier des commandes ou
d'ajouter des fonctionnalités à {\LaTeX}. On charge les paquetages
dans le préambule avec des commandes de la forme
\begin{lstlisting}
\usepackage`\marg{paquetage}'
\usepackage`\oarg{options}\marg{paquetage}'
\usepackage`\marg{paquetage1,paquetage2, ...}'
\end{lstlisting}
La première et la troisième forme permettent de charger un ou
plusieurs paquetages sans options. La seconde permet de spécifier des
\meta{options} au chargement du paquetage. Il n'est évidemment pas
possible de préciser des options avec la troisième forme puisque
{\LaTeX} ne saurait à quel paquetage celles-ci se rapportent.

Certains paquetages permettent que leurs options apparaissent parmi
les \meta{options} de la commande \cmdprint{\documentclass}. Elles
sont ainsi plus «visibles» pour d'autres paquetages. Par exemple,
l'option \code{french} que l'on retrouve dans la déclaration de la
classe à la \autoref{fig:bases:parties} est en fait une option du
paquetage \pkg{babel}.


\section{Commandes}
\label{sec:bases:commandes}

Nous avons déjà fait référence à quelques reprises au concept de
commande {\LaTeX}. Cette section se penche sur leur syntaxe.

Les formes générales des commandes {\LaTeX} sont:
\begin{lstlisting}
\`\meta{nomcommande}\oarg{arg\_optionnel}\marg{arg\_obligatoire}'
\`\meta{nomcommande}'*`\oarg{arg\_optionnel}\marg{arg\_obligatoire}'
\end{lstlisting}
Ici, \meta{nomcommande} est le nom de la commande. Il débute par le
symbole {\bs} et est exclusivement formé de lettres, habituellement
des minuscules ({\LaTeX} est sensible à la casse). La forme étoilée
d'une commande réalise généralement une action légèrement différente
de la version sans étoile. Par exemple, \cmd{\section} crée une
nouvelle section numérotée, alors que \cmd{\section*} n'insère aucune
numérotation.

Lorsque la commande prend des arguments, les arguments obligatoires
sont placés entre accolades \verb={ }= et les arguments optionnels
sont placés entre crochets \verb=[ ]=.

Certaines commandes n'ont aucun argument. Leur forme est alors
\begin{lstlisting}
\`\meta{nomcommande}
\end{lstlisting}
Dans ce cas, le nom de la commande se termine par tout symbole qui
n'est pas une lettre --- y compris l'espace! Cette règle fait en sorte
qu'une espace après le nom d'une commande est considérée comme un
marqueur de la fin du nom de la commande. Cette règle joue parfois de
vilains tours en «avalant» l'espace entre une commande et le mot qui
suit; voir l'\autoref{exemple:base:commandes} et
l'\autoref{ex:base:commandes}.

La portée d'une commande est limitée à la zone entre accolades
\verb={ }=, le cas échéant.

\begin{exemple}
  \label{exemple:base:commandes}
  Voici trois exemples de commandes {\LaTeX}: une sans argument, une
  avec un seul argument obligatoire et une commande avec deux
  arguments obligatoires et un argument optionnel.
  \begin{enumerate}
  \item La commande \cmd{\LaTeX} permet de composer de logo {\LaTeX}.
    \begin{demo}
      \begin{texample}
\begin{lstlisting}
Apprendre \LaTeX c'est
formidable!
\end{lstlisting}
        \producing
        Apprendre \LaTeX c'est
        formidable!
      \end{texample}
    \end{demo}
    On constate ici l'effet de la règle mentionnée ci-dessus:
    l'espace suivant le nom de la commande a été interprété par
    {\LaTeX} comme un marqueur de la fin de la commande et il a été
    supprimé du texte. Pour contourner ce problème, nous pouvons soit
    fournir un argument vide à la commande, soit la placer en
    accolades pour limiter sa portée à elle-même:
    \begin{demo}
      \begin{texample}
\begin{lstlisting}
Apprendre \LaTeX{} c'est
formidable!
\end{lstlisting}
        \producing
        Apprendre \LaTeX{} c'est
        formidable!
      \end{texample}
      \begin{texample}
\begin{lstlisting}
Apprendre {\LaTeX} c'est
formidable!
\end{lstlisting}
        \producing
        Apprendre {\LaTeX} c'est
        formidable!
      \end{texample}
    \end{demo}
    %
  \item La commande \cmd{\emph} met de l'emphase (en général sous
    forme d'italique) sur le ou les mots en argument.
    \begin{demo}
      \begin{texample}
\begin{lstlisting}
Il est \emph{essentiel} de
connaître la syntaxe de
{\LaTeX}.
\end{lstlisting}
        \producing
        Il est \emph{essentiel} de
        connaître la syntaxe de
        {\LaTeX}.
      \end{texample}
    \end{demo}
    %
  \item La commande \cmd{\rule} produit un rectangle plein. Elle a
    deux arguments obligatoires: la longueur et la hauteur du
    rectangle, dans l'ordre. Un argument optionnel permet de surélever
    le rectangle au-dessus de la ligne de base (voir le
    \autoref{chap:boites} pour plus de détails).
    \begin{demo}
      \begin{texample}
\begin{lstlisting}
Réglure de 1~cm de long et
3~mm d'épais surélevée de
2~points au-dessus de la
ligne de base:
\rule[2pt]{1cm}{3mm}.
\end{lstlisting}
        \producing
        Réglure de $1$~cm de long et
        $3$~mm d'épais surélevée de
        $2$~points au-dessus de la
        ligne de base: \rule[2pt]{1cm}{3mm}.
      \end{texample}
    \end{demo}
  \end{enumerate}
  \qed
\end{exemple}

\begin{exemple}
  La commande \cmd{\bfseries} sélectionne une police
  grasse pour tout le texte qui suit.
  \begin{demo}
    \begin{texample}
\begin{lstlisting}
En typographie, la \bfseries
graisse est l'épaisseur d'un
trait ou d'un caractère.
\end{lstlisting}
      \producing
      En typographie, la \bfseries
      graisse est l'épaisseur d'un
      trait ou d'un caractère.
    \end{texample}
  \end{demo}
  Pour limiter le changement à une zone de texte, il faut la délimiter
  par des accolades.
\begin{demo}
    \begin{texample}
\begin{lstlisting}
En typographie, la {\bfseries
graisse est l'épaisseur d'un
trait} ou d'un caractère.
\end{lstlisting}
      \producing
      En typographie, la {\bfseries
      graisse est l'épaisseur d’un
      trait} ou d’un caractère.
    \end{texample}
  \end{demo}
  \qed
\end{exemple}

L'utilisateur de {\LaTeX} peut définir de nouvelles commandes à
loisir. Ceci est expliqué au \autoref{chap:commandes}.


\section{Environnements}
\label{sec:bases:environnements}

Un environnement {\LaTeX} est une zone de texte délimitée par une
construction du type
\begin{lstlisting}
\begin`\marg{environnement}'
   ...
\end`\marg{environnement}'
\end{lstlisting}
Le contenu d'un environnement est traité différemment du reste du
texte en fonction des paramètres de l'environnement. Par exemple, le
texte à l'intérieur d'un environnement \Ie{center} est centré sur la
page.

Les changements induits par un environnement s'appliquent uniquement à
l'intérieur de celui-ci. Il en va de même des commandes utilisées à
l'intérieur d'un environnement.

\begin{exemple}
  L'environnement \Ie{quote} sert à composer des citations. Le texte à
  l'intérieur de l'environnement sera placé dans un bloc séparé du
  texte principal et en retrait des marges gauche et droite.
  \begin{demo}
    \begin{texample}
\begin{lstlisting}
La phrase
\begin{quote}
  Attention aux bogues dans
  le code ci-dessus; je ne
  l'ai pas testé, j'ai
  seulement prouvé qu'il
  était correct.
\end{quote}
est une citation célèbre du
créateur de {\TeX}, Donald
Knuth.
\end{lstlisting}
      \producing
      La phrase
      \begin{quote}
        Attention aux bogues dans le code ci-dessus; je ne l'ai pas testé,
        j'ai seulement prouvé qu'il était correct.
      \end{quote}
      est une citation célèbre du créateur de {\TeX}, Donald Knuth.
    \end{texample}
  \end{demo}
  Si la citation est dans la langue originale, il est préférable de la
  composer en italique.
  \begin{demo}
    \begin{texample}
\begin{lstlisting}
La phrase
\begin{quote}
  \itshape
  Beware of bugs in the above
  code; I have only proved it
  correct, not tried it.
\end{quote}
est une citation célèbre du
créateur de {\TeX}, Donald
Knuth.
\end{lstlisting}
      \producing
      La phrase
      \begin{quote}
        \itshape
        Beware of bugs in the above code; I have only proved it
        correct, not tried it.
      \end{quote}
      est une citation célèbre du créateur de {\TeX}, Donald Knuth.
    \end{texample}
  \end{demo}
  On remarque que l'effet de la commande \cmdprint{\itshape} s'est
  limité à l'environnement. %
  \qed
\end{exemple}


\section{Longueurs}
\label{sec:bases:longueurs}

Plusieurs commandes {\LaTeX} requièrent en argument une mesure de
largeur ou de hauteur. Dans la terminologie de {\LaTeX}, on parle plus
généralement de longueur\index{longueur} (\emph{length}).

Une longueur est un nombre positif, négatif ou nul
\emph{obligatoirement} et \emph{immédiatement} suivi d'un symbole
d'unité de mesure. Le \autoref{tab:bases:longueurs} présente les
principales unités de mesure utilisées par {\LaTeX} et le symbole
correspondant.

\begin{table}
  \caption{Principales unités de mesure pour les longueurs dans
    {\LaTeX}}
  \label{tab:bases:longueurs}
  \centering
  \begin{tabular}{lcl}
    \toprule
    Nom ou description & Symbole & Longueur équivalente \\
    \midrule
    millimètre & \texttt{mm} \\
    centimètre & \texttt{cm} & $10$~mm \\
    pouce      & \texttt{in} & $2,54$~cm \\
    point      & \texttt{pt} & $1/72,27$~pouce \\
    point PostScript     & \texttt{bp}   & $1/72$~pouce \\
    largeur de la lettre M & \texttt{em} & fonction de la police \\
    hauteur de la lettre x & \texttt{ex} & fonction de la police \\
    \bottomrule
  \end{tabular}
\end{table}

Il existe un certain nombre de longueurs prédéfinies. Les plus utiles
sont:
\begin{lstlisting}
\linewidth
\textwidth
\end{lstlisting}
La première contient la largeur de la ligne de texte courante et la
seconde, la largeur de la page courante. Dans du texte normal, les
deux mesures sont habituellement égales.


\section{Commentaires}
\label{sec:bases:commentaires}

Le symbole \verb=%= indique un commentaire dans le code source: tout
le texte après le symbole jusqu'à la fin de la ligne est ignoré par
{\LaTeX}.
\begin{demo}
  \begin{texample}
\begin{lstlisting}
texte % ignoré par LaTeX
\end{lstlisting}
    \producing
    texte % ignoré par LaTeX
  \end{texample}
\end{demo}


\section{Caractères spéciaux}
\label{sec:bases:caracteres}

Les claviers d'ordinateur mettent à la disposition des auteurs toutes
les lettres de l'alphabet (en versions minuscule et majuscule), les
chiffres de 0 à 9, un certain nombre de symboles et, selon le clavier,
des versions accentuées de certaines lettres. L'entrée des lettres et
des chiffres ne pose pas de problème particulier pour {\LaTeX}, mais
certains symboles sont réservés. De plus, certains symboles d'usage
courant ne sont pas disponibles sur les claviers.

\subsection{Espaces et retours à la ligne}
\label{sec:bases:caracteres:espaces}

Nous avons déjà mentionné à la \autoref{sec:bases:saisie} le
traitement spécial réservé par {\LaTeX} aux espaces et aux retours à
la ligne dans le code source. Ajoutons les précisions suivantes:
\begin{itemize}
\item seule la première espace entre deux éléments compte;
\item les espaces en début de ligne sont ignorées;
\item un retour à la ligne simple est traité comme une espace;
\item il faut deux retours à la ligne consécutifs (ce qui résulte en
  une ligne blanche dans le code source) pour identifier un changement
  de paragraphe.
\end{itemize}

Pour forcer une espace à un endroit où {\LaTeX} la supprimerait
normalement, utiliser la commande \verb*=\ = (le symbole \bs\ suivi
d'une espace, représentée ici par le symbole \verb*| |.
\begin{demo}
  \begin{texample}
\begin{lstlisting}
Apprendre \LaTeX\ c'est
formidable!
\end{lstlisting}
    \producing
    Apprendre \LaTeX\ c'est
    formidable!
  \end{texample}
\end{demo}

Dans le même ordre d'idées, le symbole \verb=~= insère une espace
insécable entre deux mots, de telle sorte que {\LaTeX} ne pourra
placer les mots sur des lignes différentes. À utiliser dans les noms ou
entre une quantité et son unité.
\begin{demo}
  \begin{texample}
\begin{lstlisting}
M.~Tremblay me doit au
moins 200~\$.
\end{lstlisting}
    \producing
    M.~Tremblay me doit au moins 200~\$.
  \end{texample}
\end{demo}

Il arrive que l'espace générée par un retour à la ligne \emph{simple}
s'avère indésirable. Dans de tels cas, placer un symbole de
commentaire \verb=%= à la fin de la ligne.
\begin{demo}
  \begin{texample}
\begin{lstlisting}
Donald Knuth est un dieu
\textsuperscript{[citation]}
.
\end{lstlisting}
    \producing
    Donald Knuth est un dieu
    \textsuperscript{[citation]}
    .
  \end{texample}
  \begin{texample}
\begin{lstlisting}
Donald Knuth est un dieu%
\textsuperscript{[citation]}%
.
\end{lstlisting}
    \producing
    Donald Knuth est un dieu%
    \textsuperscript{[citation]}%
    .
  \end{texample}
\end{demo}

\subsection{Caractères réservés}
\label{sec:bases:caracteres:reserves}

Comme à peu près tous les langages de programmation, {\TeX} réserve
certains caractères pour son usage interne. Les caractères suivants sont
interprétés comme des commandes:
\begin{center}
  \verb=# $ & ~ _ ^ % { }=
\end{center}

Pour utiliser les symboles ci-dessus tels quels dans le texte, il faut
les précéder par le symbole {\bs}:
\index{#@\cs{}\texttt{\#}}
\index{$@\cs{}\texttt{\$}} %$
\index{&@\cs{}\texttt{\&}}
\index{_@\cs{}\texttt{\_}}
\index{%@\cs{}\texttt{\%}}
\begin{demo}
  \begin{minipage}{0.15\linewidth}
    \begin{texample}
\begin{lstlisting}
\#
\end{lstlisting}
      \producing\ \#
    \end{texample}
  \end{minipage}
  \hfill
  \begin{minipage}{0.15\linewidth}
    \begin{texample}
\begin{lstlisting}
\$
\end{lstlisting}
      \producing\ \$
    \end{texample}
  \end{minipage}
  \hfill
  \begin{minipage}{0.15\linewidth}
    \begin{texample}
\begin{lstlisting}
\&
\end{lstlisting}
      \producing\ \&
    \end{texample}
  \end{minipage}
  \hfill
  \begin{minipage}{0.15\linewidth}
    \begin{texample}
\begin{lstlisting}
\_
\end{lstlisting}
      \producing\rule{0pt}{1em}\ \_
    \end{texample}
  \end{minipage}
  \\
  \begin{minipage}{0.15\linewidth}
    \begin{texample}
\begin{lstlisting}[commentstyle=\mdseries]
\%
\end{lstlisting}
      \producing\ \%
    \end{texample}
  \end{minipage}
  \hfill
  \begin{minipage}{0.15\linewidth}
    \begin{texample}
\begin{lstlisting}
\{
\end{lstlisting}
      \producing\ \{
    \end{texample}
  \end{minipage}
  \hfill
  \begin{minipage}{0.15\linewidth}
    \begin{texample}
\begin{lstlisting}
\}
\end{lstlisting}
      \producing\ \}
    \end{texample}
  \end{minipage}
  \hfill
  \begin{minipage}{0.15\linewidth}
    \mbox{}
  \end{minipage}
  \hfill
\end{demo}
(Les commandes \verb=\^= et \verb=\~= servent à créer des accents;
voir la \autoref{sec:bases:caracteres:accents}.)

\subsection{Guillemets}
\label{sec:bases:caracteres:guillemets}

L'usage des guillemets nécessite une attention particulière dans
{\LaTeX}. On ne devrait pas utiliser dans le code source les
guillemets doubles (\verb="=) que l'on trouve couramment sur les
claviers d'ordinateurs.

Pour obtenir les guillemets anglais, on utilisera deux symboles
d'accent grave côte à côte pour les guillemets ouvrants (``) et deux
apostrophes côte à côte pour les guillemets fermants ('').
\begin{demo}
  \begin{texample}
\begin{lstlisting}[escapeinside={}]
``guillemets anglais''
\end{lstlisting}
    \producing
    ``guillemets anglais''
  \end{texample}
\end{demo}

En typographie française, il convient d'utiliser les chevrons («\ »)
comme guillemets. Avec la configuration appropriée de \pkg{babel}
(voir la \autoref{sec:bases:francais}), il suffit d'utiliser dans le
code source les symboles \,\verb=«=\, et \,\verb=»=\, que l'on
retrouve normalement sur un clavier français. L'espacement requis
autour des symboles est géré automatiquement.
\begin{demo}
  \begin{texample}
\begin{lstlisting}
«guillemets français»
\end{lstlisting}
    \producing
    «guillemets français»
  \end{texample}
  \begin{texample}
\begin{lstlisting}
« guillemets français »
\end{lstlisting}
    \producing
    « guillemets français »
  \end{texample}
\end{demo}
En français, les guillemets anglais ne servent que pour les citations
incluses, c'est-à-dire pour les citations à l'intérieur d'une
citation.
\begin{demo}
  \begin{texample}
\begin{lstlisting}[escapeinside={}]
Elle m'a dit: «Laurent fait
dire ``non merci''.»
\end{lstlisting}
    \producing
    Elle m'a dit: «Laurent fait
    dire ``non merci''.»
  \end{texample}
\end{demo}

Les bons éditeurs de texte adaptés pour {\LaTeX} redéfinissent
l'action de la touche %
%% simili touche de clavier faite main à partir; utilisation de
%% \ooalign tirée de comprehensive.pdf, section 10.3
\raisebox{-2pt}{%
  \ooalign{\hfil\Large\faSquareO\hfil\cr\hfil\ttfamily\textquotedbl\hfil}} %
du clavier pour insérer les symboles appropriés selon la langue du
document.

Le paquetage \pkg{csquotes} \citep{csquotes} propose une autre approche
pour la saisie des guillemets. Il fournit une commande
\cmdprint{\enquote} qui entoure son argument des guillemets appropriés
selon le contexte et la langue du texte. Cette solution peut s'avérer
intéressante pour les nouveaux venus à {\LaTeX} qui n'ont pas encore
pris l'habitude d'entrer directement les bons guillemets.

\subsection{Trait d'union et tirets}
\label{sec:bases:caracteres:tirets}

On trouve en typographie soignée trois sortes de tirets de longueurs
différentes:
\begin{enumerate}
\item le trait d'union (-) qui sert à relier des mots entre eux;
\item le tiret demi-cadratin (--) qui sert à joindre deux éléments qui
  comportent déjà des traits d'union (\emph{Trois-Rivières--Québec});
\item le tiret cadratin (---) qui sert à identifier le changement
  d'interlocuteur dans les dialogues ou qui remplacent parfois les
  parenthèses dans du texte normal.
\end{enumerate}

S'il est simple d'entrer un trait d'union à l'ordinateur, il est rare
que les claviers comportent des touches pour entrer les tirets
demi-cadratin et cadratin (ou alors cela se fait via une obscure
combinaison de touches). {\LaTeX} permet de créer facilement tous les
tirets ci-dessus en répétant simplement le trait d'union deux ou trois
fois.
\begin{demo}
  \begin{minipage}{0.2\linewidth}
    \begin{texample}
\begin{lstlisting}
-
\end{lstlisting}
      \producing\ -
    \end{texample}
  \end{minipage}
  \hfill
  \begin{minipage}{0.2\linewidth}
    \begin{texample}
\begin{lstlisting}
--
\end{lstlisting}
      \producing\ --
    \end{texample}
  \end{minipage}
  \hfill
  \begin{minipage}{0.2\linewidth}
    \begin{texample}
\begin{lstlisting}
---
\end{lstlisting}
      \producing\ ---
    \end{texample}
  \end{minipage}
\end{demo}

\subsection{Accents et ligatures}
\label{sec:bases:caracteres:accents}

À la base, {\LaTeX} ne reconnaît pas les lettres accentuées ni les
ligatures comme {\ae} et {\oe} qui ne sont pas courantes en anglais.
Néanmoins, il est possible de créer ces symboles avec des commandes.
\begin{demo}
  \begin{minipage}{0.2\linewidth}
    \begin{texample}
\begin{lstlisting}
\'{o}
\end{lstlisting}
      \producing\ \'{o}
    \end{texample}
  \end{minipage}
  \hfill
  \begin{minipage}{0.2\linewidth}
    \begin{texample}
\begin{lstlisting}[escapeinside={}]
\`{o}
\end{lstlisting}
      \producing\ \`{o}
    \end{texample}
  \end{minipage}
  \hfill
  \begin{minipage}{0.2\linewidth}
    \begin{texample}
\begin{lstlisting}
\^{o}
\end{lstlisting}
      \producing\ \^{o}
    \end{texample}
  \end{minipage}
  \hfill
  \begin{minipage}{0.2\linewidth}
    \begin{texample}
\begin{lstlisting}
\"{o}
\end{lstlisting}
      \producing\ \"{o}
    \end{texample}
  \end{minipage} \\
  \hfill
  \begin{minipage}{0.2\linewidth}
    \begin{texample}
\begin{lstlisting}
{\ae}
\end{lstlisting}
      \producing\ \ae
    \end{texample}
  \end{minipage}
  \hfill
  \begin{minipage}{0.2\linewidth}
    \begin{texample}
\begin{lstlisting}
{\AE}
\end{lstlisting}
      \producing\ \AE
    \end{texample}
  \end{minipage}
  \hfill
  \begin{minipage}{0.2\linewidth}
    \begin{texample}
\begin{lstlisting}
{\oe}
\end{lstlisting}
      \producing\ \oe
    \end{texample}
  \end{minipage}
  \hfill
  \begin{minipage}{0.2\linewidth}
    \begin{texample}
\begin{lstlisting}
{\OE}
\end{lstlisting}
      \producing\ \OE
    \end{texample}
  \end{minipage}
\end{demo}

Évidemment, entrer tous les accents d'un texte en français avec les
commandes ci-dessus se révélerait un véritable cauchemar.
Heureusement, avec la configuration appropriée, il est possible
d'entrer directement les lettres accentuées que l'on trouve couramment
sur un clavier, voire même certaines ligatures si l'on utilise
l'encodage UTF-8. La \autoref{sec:bases:francais} fournit les détails.

Les commandes pour les accents demeurent utiles pour composer les
lettres accentuées des langues européennes autres que le français. On
trouvera la liste de ces commandes, ainsi que celles pour composer une
multitude d'autres symboles à la section~2 de la très utile %
\doc[\emph{Comprehensive {\LaTeX} Symbol List}]{comprehensive}{http://texdoc.net/pkg/comprehensive} %
\citep{comprehensive}.


\section{{\LaTeX} en français}
\label{sec:bases:francais}

Par défaut, LaTeX prévoit la rédaction de documents en anglais.
L'entrée de texte en français de façon simple et conviviale requiert
de charger quelques paquetages bien choisis. Le choix dépend lui-même
du moteur {\TeX} utilisé.

\subsection{Moteur pdf\TeX}
\label{sec:bases:francais:pdftex}

Afin de pouvoir entrer directement les lettres accentuées dans le code
source sans avoir recours aux commandes de la section précédente, il
faut s'assurer que les lettres accentuées soient reconnues par
{\LaTeX}. C'est le rôle du paquetage \pkg{inputenc} \citep{inputenc}.
Il faut spécifier le type de codage utilisé, ce qui parfois demande
quelques essais et erreurs. Le type de codage le plus universel en ce
moment est l'Unicode \citep{Unicode:5.0}. C'est le type utilisé par
défaut dans macOS et les distributions Linux; il est aussi
disponible sous Windows. Si le fichier source est sauvegardé dans le
format %
\link{https://fr.wikipedia.org/wiki/UTF-8}{UTF-8} %
de l'Unicode, on ajoutera dans le préambule la commande:
\begin{lstlisting}
\usepackage[utf8]{inputenc}
\end{lstlisting}

Les autres types de codage les plus répandus sont
l'\link{http://fr.wikipedia.org/wiki/Iso_8859-1}{ISO 8859-1}, ou
Latin-1, et l'\link{http://fr.wikipedia.org/wiki/Windows-1252}{ANSI}
de Windows. Pour utiliser ces types, on remplacera \code{utf8}
ci-dessus par, respectivement, \code{latin1} ou \code{ansiwin}.

L'entrée du texte n'est pas tout, il faut également se soucier de la
sortie, afin que les lettres accentuées apparaissent dans le document
final. {\LaTeX} utilise par défaut un codage de sortie des caractères
nommé OT1. Il est recommandé de ne pas utiliser ce codage pour les
textes en français, notamment parce que la coupure automatique ne
fonctionnera pas pour les mots contenant des lettres accentuées. On
préférera donc le plus récent codage T1, que l'on active avec le
paquetage \pkg{fontenc} \citep{fontenc}:
\begin{lstlisting}
\usepackage[T1]{fontenc}
\end{lstlisting}

La classe \class{ulthese} charge par défaut le paquetage \pkg{fontenc}
avec l'option T1 et les gabarits de documents fournis avec la classe
contiennent l'instruction pour charger \pkg{inputenc}.

%% hyperref n'aime pas \XeTeX dans un titre. Je peux utiliser un titre
%% court puisqu'il ne sera pas affiché dans la tdm (c'est une
%% sous-section).
\subsection[Moteur XeTeX]{Moteur \XeTeX}
\label{sec:bases:francais:xetex}

Le plus récent et moderne moteur {\XeTeX} utilise par défaut des
fichiers source en format Unicode UTF-8. Il n'y a donc rien de spécial
à faire pour l'entrée du texte avec ce moteur.

Pour la sortie, le paquetage \pkg{fontspec} \citep{fontspec} est
requis pour obtenir les lettres accentuées dans le fichier PDF. Cela ne
représente pas véritablement une étape de configuration additionnelle
puisque l'on charge normalement ce paquetage avec {\XeLaTeX} pour
contrôler le chargement des polices de caractères (plus de détails à
la \autoref{sec:trucs:polices}).

La classe \class{ulthese} charge par défaut le paquetage
\pkg{fontspec} avec {\XeLaTeX}.

\subsection{Typographie et mots clés français}
\label{sec:bases:francais:babel}

Entrer du texte en français dans le code source n'est pas tout. Il
faut adapter {\LaTeX} au français, qu'il s'agisse des mots clés (tels
«Table des matières» ou «Bibliographie»), de la typographie ou de la
coupure des mots. La solution standard provient du paquetage
\pkg{babel} \citep{babel}. Celui-ci permet de combiner plusieurs
langues dans un même document et de passer de l'une à l'autre
facilement. Consulter la %
\doc{frenchb}{http://texdoc.net/pkg/babel-french/} %
pour connaître toutes les adaptations au français offertes par ce
paquetage incontournable.

Le paquetage \pkg{babel} est chargé par défaut par la classe
\class{ulthese}.

\subsection{Guillemets français}
\label{sec:bases:francais:guillemets}

Une espace fine devrait séparer les guillemets du ou des mots qu'ils
entourent. Pour que l'espace autour des symboles soit géré
automatiquement par \pkg{babel}, il suffit d'ajouter dans le préambule
la commande
\begin{lstlisting}
\frenchbsetup{og=«, fg=»}
\end{lstlisting}
On trouve cette configuration dans les gabarits livrés avec la classe
\class{ulthese}.

\subsection{Séparateur décimal}
\label{sec:bases:francais:virgule}

Le séparateur décimal en français est la virgule et non le point. Or,
en mode mathématique, LaTeX ajoute automatiquement une espace fine
après une virgule comme s'il s'agissait d'une énumération.

Afin de pouvoir utiliser de manière conviviale la virgule comme
séparateur décimal en mode mathématique, on chargera le très pratique
paquetage \pkg{icomma} \citep{icomma}. Avec ce paquetage, la virgule
agira comme séparateur décimal en mode mathématique seulement
lorsqu'elle est suivie d'un caractère autre que l'espace.
\begin{demo}
  \begin{minipage}{0.5\linewidth}
    \begin{texample}[0.6\linewidth]
\begin{lstlisting}
% nombre décimal
$x = 1,2$
\end{lstlisting}
      \producing
      $x = 1,2$
    \end{texample}
  \end{minipage}
  \hfill
  \begin{minipage}{0.5\linewidth}
    \begin{texample}[0.6\linewidth]
\begin{lstlisting}
% énumération
$x = 1, 2, 3$
\end{lstlisting}
      \producing
      $x = 1, 2, 3$
    \end{texample}
  \end{minipage}
\end{demo}

\begin{important}
  Le paquetage \pkg{icomma} peut interagir avec d'autres. Aussi
  est-il préférable de le charger parmi les derniers dans le préambule, et
  certainement après les paquetages \pkg{fontspec} et
  \pkg{unicode-math} avec {\XeLaTeX}.
\end{important}

\subsection{Séparateur des milliers}
\label{sec:bases:francais:milliers}

Le séparateur des milliers en français est l'espace. Lorsque le
paquetage \pkg{numprint} \citep{numprint} est chargé, \pkg{babel}
fournit la commande \cmd{\nombre} pour formater automatiquement les
nombres.
\begin{demo}
  \begin{texample}
\begin{lstlisting}
\nombre{123456789}
\end{lstlisting}
    \producing
    \nombre{123456789}
  \end{texample}
\end{demo}

Le \autoref{tab:bases:francais} résume les divers enjeux liés à la
rédaction en français avec {\LaTeX} et les paquetages qui offrent des
solutions.

\begin{table}
  \centering
  \caption{Enjeux de rédaction en français et paquetages offrant des
    solutions}
  \label{tab:bases:francais}
  \begin{tabularx}{1.0\linewidth}{Xl}
    \toprule
    Enjeu & Solution \\
    \midrule
    \addlinespace[0.5\normalbaselineskip]
    Traduction des mots clés prédéfinis & \pkg{babel} \\
    \addlinespace[0.5\normalbaselineskip]
    Coupure de mots & \pkg{babel} \\
    \addlinespace[0.5\normalbaselineskip]
    Typographie française & \pkg{babel} \\
    \addlinespace[0.5\normalbaselineskip]
    Lettres accentuées dans source & \pkg{inputenc} (\LaTeX) \\
                                   & source en UTF-8 (\XeLaTeX) \\
    \addlinespace[0.5\normalbaselineskip]
    Lettres accentuées dans sortie & \pkg{fontenc} (\LaTeX) \\
                                   & \pkg{fontspec} (\XeLaTeX) \\
    \addlinespace[0.5\normalbaselineskip]
    Virgule comme séparateur décimal & \pkg{icomma} \\
    \addlinespace[0.5\normalbaselineskip]
    Espace comme séparateur des milliers & \pkg{numprint} et
                                           \pkg{babel} \\
    \bottomrule
  \end{tabularx}
\end{table}


%%%
%%% Exercices
%%%

\section{Exercices}
\label{sec:bases:exercices}

\Opensolutionfile{solutions}[solutions-bases]

\begin{Filesave}{solutions}
\section*{Chapitre \ref*{chap:bases}}
\addcontentsline{toc}{section}{Chapitre \protect\ref*{chap:bases}}

\end{Filesave}

\begin{exercice}[nosol]
  Utiliser le fichier \fichier{exercice\_minimal.tex}.
  \begin{enumerate}
  \item Compiler le document avec la classe \class{article}, puis avec
    la classe \class{book}. Observer le résultat.
  \item Ajouter du texte en français (avec accents), puis compiler et
    observer le résultat.
  \end{enumerate}
\end{exercice}

\begin{exercice}
  \label{ex:base:commandes}
  Modifier le fichier \fichier{exercice\_commandes.tex} afin
  de produire le texte ci-dessous.
  \begin{demo}
    \fbox{\includegraphics[viewport=108 551 502 665,clip=true,width=0.98\linewidth]{exercice_commandes-solution}}
  \end{demo}
  \begin{sol}
    Il y a trois modifications principales à apporter au fichier.
    \begin{enumerate}[1.]
    \item La commande {\LaTeX} qui imprime le logo de LaTeX n'ayant
      pas d'argument, il faut l'écrire sous la forme
      \verb={\LaTeX}= ou \verb=\LaTeX{}= pour éviter que l'espace
      qui suit ne soit avalée par la commande.
    \item Il faut délimiter par des accolades la zone à laquelle la
      commande \cmdprint{\bfseries} doit s'appliquer:
\begin{lstlisting}
{\bfseries limitée}
\end{lstlisting}
    \item Il faut utiliser l'environnement \Pe{enumerate} pour les
      listes numérotées. L'environnement \Pe{itemize} sert pour les
      listes à puces.
    \end{enumerate}
  \end{sol}
\end{exercice}

\begin{exercice}[nosol]
  \begin{enumerate}
  \item Compiler le fichier \fichier{exercice\_classe+paquetages.tex}.
  \item Changer la police du document pour 11~points, puis 12~points.
    Observer l'effet sur les marges et sur la coupure automatique des
    mots dans le document compilé.
  \item Activer le paquetage \pkg{icomma} en supprimant le symbole
    \% au début de la ligne dans le préambule. Compiler le document
    et observer l'effet sur la formule mathématique.
  \item Charger le paquetage \pkg{numprint} avec l'option
    \verb=autolanguage= (\emph{après} le paquetage \pkg{babel}). Dans
    le code source de la formule mathématique, changer
\begin{lstlisting}
10 000
\end{lstlisting}
    pour
\begin{lstlisting}
\nombre{10000}
\end{lstlisting}
    Compiler et observer le résultat.
  \end{enumerate}
\end{exercice}

\Closesolutionfile{solutions}

%%% Local Variables:
%%% mode: latex
%%% TeX-engine: xetex
%%% TeX-master: "formation-latex-ul"
%%% coding: utf-8
%%% End:

version https://git-lfs.github.com/spec/v1
oid sha256:18d6ce7550cc3d788899ad02ead226c3f21159327b45f367d344ff5e4536ce21
size 38391

version https://git-lfs.github.com/spec/v1
oid sha256:63bdd8292fea5691521a32972fd47b9b9e58fe984b40afc1b7086462f5a5f5bd
size 19777

version https://git-lfs.github.com/spec/v1
oid sha256:0854b2c07bc12e8c8431853876afbfb9f1fd81fbc4328ca55362dc5e6bf55291
size 20184

version https://git-lfs.github.com/spec/v1
oid sha256:f0de4cc96605f95da458c00f64cde1e2521258ced0097498f9fc358191bdfcb5
size 43788

\chapter{Mathématiques}
\label{chap:math}


S'il est un domaine où {\LaTeX} brille particulièrement, c'est bien
dans la préparation et la présentation d'équations mathématiques ---
des plus simples aux plus complexes. Après tout, l'amélioration de la
qualité typographique des équations mathématiques dans son ouvrage
phare \emph{The Art of Computer Programming} figurait parmi les
objectifs premiers de Knuth lorsqu'il a développé {\TeX}.


\section{Principes de base du mode mathématique}
\label{sec:math:base}

La mise en forme d'équations mathématiques requiert d'indiquer à
l'ordinateur, dans un langage spécial, le contenu des dites équations
et la position des symboles: en exposant, en indice, sous forme de
fraction, etc. L'ordinateur peut ensuite assembler le tout à partir de
règles typographiques portant, par exemple, sur la représentation des
variables et des constantes, l'espacement entre les symboles ou la
disposition des équations selon qu'elles apparaissent au fil du texte
ou hors d'un paragraphe.

On indique à {\LaTeX} que l'on change de «langage», par l'utilisation
d'un mode mathématique. Il y a deux grandes manière d'activer le mode
mathématique:
\begin{enumerate}
\item en insérant le code entre les symboles \verb=$ $= pour générer
  une équation «en ligne», ou au fil du texte;
  \begin{demo}
    \begin{texample}
\begin{lstlisting}
on sait que $(a + b)^2 =
a^2 + 2ab + b^2$, d'où
on obtient...
\end{lstlisting}
      \producing
      on sait que $(a + b)^2 = a^2 + 2ab + b^2,$ d'où on obtient...
    \end{texample}
  \end{demo}
\item en utilisant un environnement servant à créer une équation hors
  paragraphe;
  \begin{demo}
    \begin{texample}
\begin{lstlisting}
on sait que
\begin{equation*}
  (a + b)^2
  = a^2 + 2ab + b^2,
\end{equation*}
d'où on obtient...
\end{lstlisting}
      \producing
      on sait que
      \begin{equation*}
        (a + b)^2 = a^2 + 2ab + b^2,
      \end{equation*}
      d'où on obtient...
    \end{texample}
  \end{demo}
\end{enumerate}

Dans l'exemple ci-dessus, l'environnement \Pe{equation*} (tiré du
paquetage \pkg{amsmath}, voir la section suivante) crée une
équation hors paragraphe, centrée sur la ligne et non numérotée. Avec
l'environnement \Ie{equation} (donc sans \verb=*= dans le nom),
{\LaTeX} ajoute automatiquement un numéro d'équation séquentiel aligné
sur la marge de droite:
\begin{demo}
  \begin{texample}
\begin{lstlisting}
on sait que
\begin{equation}
  (a + b)^2 = a^2 + b^2,
\end{equation}
d'où on obtient...
\end{lstlisting}
    \producing
    on sait que
    \begin{equation}
      (a + b)^2 = a^2 + b^2,
    \end{equation}
    d'où on obtient...
  \end{texample}
\end{demo}
Cette disposition est la plus usuelle dans les ouvrages mathématiques.
Le type de numérotation diffère selon qu'un document comporte des
chapitres ou non.

En mode mathématique, les chiffres sont automatiquement considérés
comme des constantes, les lettres comme des variables et une suite de
lettres comme un produit de variables (nous verrons plus loin comment
représenter des fonctions mathématiques comme $\sin$, $\log$ ou
$\lim$). Ceci a trois conséquences principales:
\begin{enumerate}
\item conformément aux conventions typographiques, les chiffres sont
  représentés en caractère \textrm{romain} et les variables, en
  italique;
  \begin{demo}
    \begin{texample}
\begin{lstlisting}
$123xyz$
\end{lstlisting}
      \producing
      $123xyz$
    \end{texample}
  \end{demo}
\item l'espace entre les constantes, les variables et les opérateurs
  mathématiques est géré automatiquement;
  \begin{demo}
    \begin{texample}
\begin{lstlisting}
$z = 2 x + 3 x y$
\end{lstlisting}
    \producing
    $z = 2 x + 3 x y$
    \end{texample}
  \end{demo}
\item les espaces dans le code source n'ont aucun impact sur la
  disposition d'une équation.
  \begin{demo}
    \begin{texample}
\begin{lstlisting}
$z=2x + 3xy$
\end{lstlisting}
    \producing
    $z=2x + 3xy$
    \end{texample}
  \end{demo}
\end{enumerate}

Quant au langage retenu par {\LaTeX} pour décrire les équations
mathématiques, il est très similaire à celui que l'on utiliserait pour
le faire à voix haute. Il faut simplement recourir à des commandes
pour identifier les symboles mathématiques que l'on ne retrouve pas
sur un clavier usuel, comme les lettres grecques, les opérateurs
d'inégalité ou les symboles de sommes et d'intégrales.


\section{Un paquetage incontournable}
\label{sec:math:amsmath}

Le paquetage \pkg{amsmath} \citep{amsmath} produit par la prestigieuse
\emph{American Mathematical Society} fournit diverses extensions à
{\LaTeX} pour faciliter encore davantage la saisie d'équations
mathématiques complexes et en améliorer la présentation. L'utilisation
de ce paquetage doit être considérée incontournable pour tout document
contenant plus que quelques équations très simples.

Au chapitre des améliorations fournies par \pkg{amsmath}, notons
particulièrement:
\begin{itemize}
\item plusieurs environnements pour les équations hors paragraphe, en
  particulier pour les équations multilignes;
\item une meilleure gestion de l'espacement autour des symboles de
  relation (comme les signes d'égalité) dans les équations
  multilignes;
\item une commande pour faciliter l'entrée de texte à l'intérieur du
  mode mathématique;
\item un environnement pour la saisie des matrices et des coefficients
  binomiaux;
\item des commandes pour les intégrales multiples;
\item la possibilité de définir de nouveaux opérateurs mathématiques.
\end{itemize}
Nous décrivons certaines de ces fonctionnalités dans la suite, mais
tout utilisateur du paquetage devrait impérativement consulter sa %
\doc[documentation complète]{amsldoc}{http://texdoc.net/pkg/amsmath}.


\section{Principaux éléments du mode mathématique}
\label{sec:math:bases}

Cette section explique comment créer et assembler divers éléments
d'une formule mathématique: exposants, indices, fractions, texte, etc.
Les seuls symboles utilisés sont pour le moment les chiffres et les
lettres latines. La présentation d'une partie de l'éventail de
symboles mathématiques offerts par {\LaTeX} fera l'objet de la
\autoref{sec:math:symboles}.

\subsection{Exposants et indices}
\label{sec:math:bases:exposants}

{\LaTeX} permet de créer facilement et avec la bonne taille de
symboles n'importe quelle combinaison d'exposants et d'indices.

On place un caractère en \textsuperscript{exposant} avec la
commande \verb=^= et en \textsubscript{indice} avec la commande
\verb=_=. Les indices et exposants se combinent naturellement.
\begin{demo}
  \def\strut{\rule[-0.4ex]{0pt}{2ex}}
  \begin{minipage}{0.3\linewidth}
    \begin{texample}[0.6\linewidth]
\begin{lstlisting}
x^2
\end{lstlisting}
      \producing\strut $x^2$
    \end{texample}
  \end{minipage}
  \quad
  \begin{minipage}{0.3\linewidth}
    \begin{texample}[0.6\linewidth]
\begin{lstlisting}
a_n
\end{lstlisting}
      \producing\strut $a_n$
    \end{texample}
  \end{minipage}
  \quad
  \begin{minipage}{0.3\linewidth}
    \begin{texample}[0.63\linewidth]
\begin{lstlisting}
x_i^k
\end{lstlisting}
      \producing\strut $x_i^k$
    \end{texample}
  \end{minipage}
\end{demo}
(L'ordre de saisie n'a pas d'importance; le troisième exemple
donnerait le même résultat avec \verb=x^k_i=.)

Si l'exposant ou l'indice compte plus d'un caractère, il faut
regrouper le tout entre accolades \verb={ }=.
\begin{demo}
  \begin{minipage}{0.3\linewidth}
    \begin{texample}[0.6\linewidth]
\begin{lstlisting}
x^{2k + 1}
\end{lstlisting}
      \producing\strut $x^{2k + 1}$
    \end{texample}
  \end{minipage}
  \quad
  \begin{minipage}{0.3\linewidth}
    \begin{texample}[0.6\linewidth]
\begin{lstlisting}
x_{i,j}
\end{lstlisting}
      \producing\strut $x_{i,j}$
    \end{texample}
  \end{minipage}
  \quad
  \begin{minipage}{0.3\linewidth}
    \begin{texample}[0.63\linewidth]
\begin{lstlisting}
x_{ij}^{2n}
\end{lstlisting}
      \producing\strut $x_{ij}^{2n}$
    \end{texample}
  \end{minipage}
\end{demo}

Toutes les combinaisons d'exposants et d'indices sont possibles, y
compris les puissances de puissances ou les indices d'indices.
\begin{demo}
  \begin{minipage}{0.3\linewidth}
    \begin{texample}[0.6\linewidth]
\begin{lstlisting}
e^{-x^2}
\end{lstlisting}
      \producing\strut $e^{-x^2}$
    \end{texample}
  \end{minipage}
  \quad
  \begin{minipage}{0.6\linewidth}
    \begin{texample}[0.6\linewidth]
\begin{lstlisting}
A_{i_s, k^n}^{y_i}
\end{lstlisting}
      \producing\strut $A_{i_s,k^n}^{y_i}$
    \end{texample}
  \end{minipage}
  \quad \mbox{} % pour alignement avec bloc d'exemples ci-dessus
\end{demo}

\begin{important}
  Les commandes \verb=^= et \verb=_= sont permises dans le mode
  mathématique seulement. En fait, si {\TeX} rencontre l'une de ces
  commandes en mode texte, il tentera automatiquement de passer au
  mode mathématique après avoir émis l'avertissement
\begin{verbatim}
! Missing $ inserted.
\end{verbatim}
  Il est assez rare que le résultat soit celui souhaité.
\end{important}

\subsection{Fractions}
\label{sec:math:bases:fractions}

Il y a plusieurs façons de représenter une fraction selon qu'elle se
trouve au fil du texte, dans une équation hors paragraphe ou à
l'intérieur d'une autre fraction.

Pour les fractions au fil du texte, il vaut souvent mieux utiliser
simplement la barre oblique \verb=/= pour séparer le numérateur du
dénominateur, quitte à utiliser des parenthèses. Ainsi, on utilise
\verb=$(n + 1)/2$= pour obtenir $(n + 1)/2$.

De manière plus générale, la commande
\begin{lstlisting}
\frac`\marg{numérateur}\marg{dénominateur}'
\end{lstlisting}
dispose \meta{numérateur} au-dessus de \meta{dénominateur}, séparé par
une ligne horizontale. La taille des caractères s'ajuste
automatiquement selon que la fraction se trouve au fil du texte ou
dans une équation hors paragraphe, ainsi que selon la position de la
fraction dans l'équation.
\begin{demo}
  \begin{texample}
\begin{lstlisting}
% taille au fil du texte
On a $z_1 = \frac{x}{y}$ et
$z_2 = xy$.
\end{lstlisting}
    \producing On a $z_1 = \frac{x}{y}$ et $z_2 = xy$.
  \end{texample}

  \begin{texample}
\begin{lstlisting}
% taille hors paragraphe
On a
\begin{equation*}
  z_1 = \frac{x}{y}
\end{equation*}
et $z_2 = xy$.
\end{lstlisting}
    \producing On a
    \begin{equation*}
      z_1 = \frac{x}{y}
    \end{equation*}
    et $z_2 = xy$.
  \end{texample}

  \begin{texample}
\begin{lstlisting}
% deux tailles combinées
Soit
\begin{equation*}
  z = \frac{\frac{x}{2}
    + 1}{y}.
\end{equation*}
\end{lstlisting}
    \producing Soit
    \begin{equation*}
      z = \frac{\frac{x}{2} + 1}{y}.
    \end{equation*}
  \end{texample}
\end{demo}

Les commandes
\begin{lstlisting}
\dfrac`\marg{numérateur}\marg{dénominateur}'
\tfrac`\marg{numérateur}\marg{dénominateur}'
\end{lstlisting}
de \pkg{amsmath} permettent de forcer une fraction à adopter la taille
d'une fraction hors paragraphe (\emph{displayed}) dans le cas de
\cmd{\dfrac}, ou de celle d'une fraction au fil du texte (\emph{text})
dans le cas de \cmd{\tfrac}. Consulter l'\autoref{ex:math:matrices} à
la \autopageref{ex:math:matrices} pour visualiser l'effet de la
commande \cmd{\dfrac}.

\begin{conseil}
  Il est parfois visuellement plus intéressant, surtout au fil du
  texte, d'écrire une fraction comme $1/x$ sous la forme $x^{-1}$.
\end{conseil}

\subsection{Racines}
\label{sec:math:bases:racines}

La commande
\begin{lstlisting}
\sqrt`\oarg{n}\marg{radicande}'
\end{lstlisting}
construit un symbole de radical autour de \meta{radicande}, par défaut
la racine carrée. Si l'argument optionnel \meta{n} est spécifié, c'est
plutôt un symbole de racine d'ordre $n$ qui est tracé. La longueur et
la hauteur du radical s'adapte toujours à celles du radicande.
\begin{demo}
  \begin{minipage}{0.3\linewidth}
    \begin{texample}[0.6\linewidth]
\begin{lstlisting}
\sqrt{2}
\end{lstlisting}
      \producing $\sqrt{2}$
    \end{texample}
  \end{minipage}
  \quad
  \begin{minipage}{0.3\linewidth}
    \begin{texample}[0.6\linewidth]
\begin{lstlisting}
\sqrt{625}
\end{lstlisting}
      \producing
      $\sqrt{625}$
    \end{texample}
  \end{minipage}
  \quad
  \begin{minipage}{0.3\linewidth}
    \begin{texample}[0.6\linewidth]
\begin{lstlisting}
\sqrt[3]{8}
\end{lstlisting}
      \producing
      $\sqrt[3]{8}$
    \end{texample}
  \end{minipage}

  \begin{texample}
\begin{lstlisting}
\sqrt[n]{x + y + z}
\end{lstlisting}
    \producing
    $\sqrt[n]{x + y + z}$
  \end{texample}

  \begin{texample}
\begin{lstlisting}
\sqrt{\frac{x + y}{x^2 - y^2}}
\end{lstlisting}
    \producing
    $\displaystyle \sqrt{\frac{x + y}{x^2 - y^2}}$
  \end{texample}
\end{demo}

\subsection{Sommes et intégrales}
\label{sec:math:bases:sommes-et-integrales}

Les sommes et les intégrales requièrent un symbole spécial ainsi que
des limites inférieures et supérieures, le cas échéant.

Les commandes \cmd{\sum} et \cmd{\int} servent respectivement à tracer
les symboles de somme $\sum$ et d'intégrale $\int$. Le paquetage
\pkg{amsmath} fournit également des commandes comme \cmd{\iint} et
\cmd{\iiint} pour obtenir des symboles d'intégrales multiples finement
disposés ($\iint$ et $\iiint$).

On entre les éventuelles limites inférieures et supérieures comme des
indices et des exposants.
\begin{demo}
  \begin{texample}
\begin{lstlisting}
\sum_{i = 0}^n x_i
\end{lstlisting}
    \producing $\displaystyle \sum_{i = 0}^n x_i$
  \end{texample}

  \begin{texample}
\begin{lstlisting}
\int_0^{10} f(x)\, dx
\end{lstlisting}
    \producing $\displaystyle \int_0^{10} f(x)\, dx$
  \end{texample}

  \begin{texample}
\begin{lstlisting}
\iint_D f(x, y)\, dx dy
\end{lstlisting}
    \producing $\displaystyle \iint_D f(x, y)\, dx dy$
  \end{texample}
\end{demo}

La taille des symboles et la position des limites s'ajustent
automatiquement selon le contexte. Au fil du texte, la somme et
l'intégrale simple ci-dessus apparaîtraient comme $\sum_{i = 0}^n x_i$
et $\int_0^{10} f(x)\, dx$.

Dans une intégrale il est recommandé de séparer l'intégrande de
l'opérateur de différentiation $dx$ par une espace fine. C'est ce à
quoi sert la commande \cmd{\,} ci-dessus; voir aussi le
\autoref{tab:math:espaces} de la \autopageref{tab:math:espaces}.

\subsection{Points de suspension}
\label{sec:math:bases:dots}

Les formules mathématiques comportent fréquemment des points de
suspension dans des suites de variables ou d'opérations. On recommande
d'éviter de les entrer comme trois points finaux consécutifs, car
l'espacement entre les points sera trop petit et le résultat, jugé
disgracieux\footnote{%
  Le résultat exact dépend de la police utilisée.} %
d'un point de vue typographique: $...$

Le \autoref{tab:math:dots} fournit les commandes {\LaTeX} servant à
générer divers types de points de suspension.

\begin{table}
  \caption{Points de suspension}
  \label{tab:math:dots}
  \centering
  \begin{tabular}{lll}
    \toprule
    commande & type de points & exemple \\
    \midrule
    \cmd{\dots} &  sélection automatique \\
    \cmd{\ldots} & points à la ligne de base & $x_1, \ldots, x_n$ \\
    \cmd{\cdots} & points centrés & $x_1 + \cdots + x_n$ \\
    \addlinespace[0.5\normalbaselineskip]
    \cmd{\vdots} & points verticaux & $
                                      \begin{matrix}
                                        x_1 \\ \vdots \\ x_n
                                      \end{matrix}$ \\
    \addlinespace[0.5\normalbaselineskip]
    \cmd{\ddots} & points diagonaux & $
                                      \begin{matrix}
                                        x_1 &&\\ &\ddots& \\ && x_n
                                      \end{matrix}$ \\
    \bottomrule
  \end{tabular}
\end{table}

Avec \pkg{amsmath}, la commande \cmd{\dots} tâche de sélectionner
automatiquement entre les points à la ligne de base ou les points
centrés selon le contexte. Comme le résultat est en général le bon,
nous recommandons d'utiliser principalement cette commande pour
insérer des points de suspension en mode mathématique.
\begin{demo}
  \begin{texample}
\begin{lstlisting}
$x_1, \dots, x_n$
\end{lstlisting}
    \producing $x_1, \dots, x_n$
  \end{texample}

  \begin{texample}
\begin{lstlisting}
$x_1, \ldots, x_n$
\end{lstlisting}
    \producing $x_1, \ldots, x_n$
  \end{texample}

  \begin{texample}
\begin{lstlisting}
$x_1 + \dots + x_n$
\end{lstlisting}
    \producing $x_1 + \dots + x_n$
  \end{texample}

  \begin{texample}
\begin{lstlisting}
$x_1 + \cdots + x_n$
\end{lstlisting}
    \producing $x_1 + \cdots + x_n$
  \end{texample}
\end{demo}
Le paquetage définit également les commandes sémantiques suivantes:
\begin{itemize}
\item \cmd{\dotsc} pour des «points avec des virgules» (\emph{commas});
\item \cmd{\dotsb} pour des «points avec des opérateurs binaires»;
\item \cmd{\dotsm} pour des «points de multiplication»;
\item \cmd{\dotsi} pour des «points avec des intégrales»;
\item \cmd{\dotso} pour d'«autres points» (\emph{other}).
\end{itemize}

\subsection{Texte et espaces}
\label{sec:math:bases:texte}

On l'a vu, en mode mathématique {\LaTeX} traite les lettres comme des
variables et gère automatiquement l'espacement entre les divers
symboles. Or, il n'est pas rare que des formules mathématiques
contiennent du texte (notamment des mots comme «où», «si», «quand»).
De plus, il est parfois souhaitable de pouvoir ajuster les blancs
entre des éléments.

La commande de \pkg{amsmath}
\begin{lstlisting}
\text`\marg{texte}'
\end{lstlisting}
insère \meta{texte} dans une formule mathématique. Le texte est inséré
tel quel, sans aucune gestion des espaces avant ou après le texte. Si
des espaces sont nécessaires, ils doivent faire partie de
\meta{texte}.
\begin{demo}
  \begin{texample}
\begin{lstlisting}
f(x) = a e^{-ax}
\text{ pour } x > 0
\end{lstlisting}
    \producing $f(x) = a e^{-ax} \text{ pour } x > 0$
  \end{texample}
\end{demo}

Les commandes
\begin{lstlisting}
\quad
\qquad
\end{lstlisting}
insèrent un blanc de largeur variable selon la taille de la police en
vigueur. La commande \cmd{\quad} insère un blanc de $1$~em (la largeur de
la lettre M dans la police en vigueur), alors que \cmd{\qquad}
insère le double de cette longueur\footnote{%
  Bien qu'elles soient surtout utilisées dans le mode mathématique,
  les commandes \cmd{\quad} et \cmd{\qquad} sont également valides
  dans le mode texte.}.
\begin{demo}
  \begin{texample}
\begin{lstlisting}
f(x) = a e^{-ax},
\quad x > 0
\end{lstlisting}
    \producing
    $f(x) = a e^{-ax}, \quad x > 0$
  \end{texample}
\end{demo}

Le \autoref{tab:math:espaces} répertorie et compare les différentes
commandes qui permettent d'insérer des espaces plus ou moins fines
entre des éléments dans le mode mathématique.

\begin{table}
  \caption{Espaces dans le mode mathématique}
  \label{tab:math:espaces}
  \centering
  \begin{tabular}{lll}
    \toprule
    Commande & Longueur & Exemple \\
    \midrule
             & pas d'espace & \spx{} \\
    \cmd{\,} & $3/18$ de \cmdprint{quad} & \spx{\,} \\
    \cmd{\:} & $4/18$ de \cmdprint{quad} & \spx{\:} \\
    \cmd{\;} & $5/18$ de \cmdprint{quad} & \spx{\;} \\
    \pixabang & $-3/18$ de \cmdprint{quad} & \spx{\!} \\
    \cmd{\quad} & $1$~em & \spx{\quad} \\
    \cmd{\qquad} & $2$~em & \spx{\qquad} \\
    \bottomrule
  \end{tabular}
\end{table}

\subsection{Fonctions et opérateurs}
\label{sec:math:bases:fonctions}

Les règles de typographie pour les équations mathématiques exigent que
les variables apparaissent en \textit{italique}, mais que les noms de
fonctions, eux, apparaissent en \textrm{romain}, comme le texte
standard. Pensons, ici, à des fonctions comme $\sin$ ou $\log$.

On sait que {\LaTeX} interprétera un nom de fonction saisi tel quel
comme un produit de variables:
\begin{demo}
  \begin{minipage}{0.45\linewidth}
    \begin{texample}
\begin{lstlisting}
$2 sin 30$
\end{lstlisting}
      \producing
      $2 sin 30$
    \end{texample}
  \end{minipage}
  \hfill
  \begin{minipage}{0.45\linewidth}
    \begin{texample}
\begin{lstlisting}
$3 log 2$
\end{lstlisting}
      \producing
      $3 log 2$
    \end{texample}
  \end{minipage}
\end{demo}
La commande \cmdprint{\text} permettrait d'afficher les noms de
fonction en romain, mais l'utiliser à répétition se révélerait
rapidement peu pratique. De plus, un nom de fonction n'est pas du
texte à proprement parler, mais plutôt un opérateur mathématique.

La solution de {\LaTeX} consiste à définir des commandes pour un grand
nombre de fonctions et d'opérateurs mathématiques standards:
\begin{lstlisting}
\arccos   \cosh    \det    \inf      \limsup   \Pr     \tan
\arcsin   \cot     \dim    \ker      \ln       \sec    \tanh
\arctan   \coth    \exp    \lg       \log      \sin
\arg      \csc     \gcd    \lim      \max      \sinh
\cos      \deg     \hom    \liminf   \min      \sup
\end{lstlisting}
L'espacement autour des fonctions et opérateurs est géré par {\LaTeX}.
\begin{demo}
  \begin{minipage}{0.45\linewidth}
    \begin{texample}
\begin{lstlisting}
$2 \sin30$
\end{lstlisting}
      \producing
      $2 \sin30$
    \end{texample}
  \end{minipage}
  \hfill
  \begin{minipage}{0.45\linewidth}
    \begin{texample}
\begin{lstlisting}
$3 \log 2$
\end{lstlisting}
      \producing
      $3 \log 2$
    \end{texample}
  \end{minipage}
\end{demo}

Certaines des fonctions ci-dessus, notamment \cmd{\lim}, acceptent des
limites comme les symboles de somme et d'intégrale.
\begin{demo}
  \begin{texample}
\begin{lstlisting}
% au fil du texte
\lim_{x \to 0} x = 0
\end{lstlisting}
    \producing $\lim_{x \to 0} x = 0$
  \end{texample}

  \begin{texample}
\begin{lstlisting}
% hors paragraphe
\lim_{x \to 0} x = 0
\end{lstlisting}
    \producing $\displaystyle \lim_{x \to 0} x = 0$
  \end{texample}
\end{demo}

La commande \cmd{\DeclareMathOperator} de \pkg{amsmath} permet de
définir de nouveaux opérateurs mathématiques lorsque nécessaire;
consulter la %
\doc{amsldoc}{http://texdoc.net/pkg/amsmath} %
du paquetage (section~5.1) pour les détails.


\begin{exemple}
  Le matériel passé en revue jusqu'à maintenant permet déjà
  de composer des équations élaborées --- sous réserve qu'elles
  tiennent sur une seule ligne comme dans le présent exemple.

  On présente ci-dessous, pièce par pièce, le code {\LaTeX} pour créer
  l'équation suivante:
  \begin{equation*}
    \int_x^\infty (y - x) f_{X|X > x}(y)\, dy =
    \frac{1}{1 - F_X(x)} \int_x^\infty (y - x) f_X(y)\, dy.
  \end{equation*}
  \begin{demo}
    \begin{texample}[0.58\linewidth]
\begin{lstlisting}
\begin{equation*}
\end{lstlisting}
      \producing
      \emph{équation hors paragraphe}
    \end{texample}

    \begin{texample}[0.58\linewidth]
\begin{lstlisting}
\int_x^\infty
\end{lstlisting}
      \producing
      $\displaystyle \int_x^\infty$
    \end{texample}

    \begin{texample}[0.58\linewidth]
\begin{lstlisting}
(y - x) f_{X|X > x}(y)\, dy =
\end{lstlisting}
      \producing
      $(y - x) f_{X|X > x}(y)\, dy =$
    \end{texample}

    \begin{texample}[0.58\linewidth]
\begin{lstlisting}
\frac{1}{1 - F_X(x)}
\end{lstlisting}
      \producing
      $\dfrac{1}{1 - F_X(x)}$
    \end{texample}

    \begin{texample}[0.58\linewidth]
\begin{lstlisting}
\int_x^\infty (y - x) f_X(y)\, dy
\end{lstlisting}
      \producing
      $\displaystyle \int_x^\infty (y - x) f_X(y)\, dy$
    \end{texample}

    \begin{texample}[0.58\linewidth]
\begin{lstlisting}
\end{equation*}
\end{lstlisting}
      \producing
      \emph{fin de l'environnement}
    \end{texample}
  \end{demo}
  \qed
\end{exemple}


\section{Symboles mathématiques}
\label{sec:math:symboles}

Outre les chiffres et les lettres de l'alphabet, les claviers
d'ordinateurs ne comptent normalement que les symboles
mathématiques suivants:
\begin{center}
%\begin{verbatim}
\verb0+  -  =  <  >  /  :  !  '  |  [  ]  (  ) { }0
%\end{verbatim}
\end{center}
Ils s'emploient directement dans les équations, sauf les accolades
\verb={ }= qui sont des symboles réservés par {\LaTeX}. Il faut donc
entrer celles-ci avec \cmdprint{\{} et \cmdprint{\}}, comme dans du
texte normal.

Pour représenter les innombrables autres symboles mathématiques, on
aura recours à des commandes qui débutent, comme d'habitude, par le
symbole {\bs} et dont le nom est habituellement dérivé de la
signification mathématique du symbole.

Si un symbole mathématique a été utilisé quelque part dans une
publication, il y a de fortes chances que sa version existe dans
{\LaTeX}. Il serait donc utopique de tenter de faire ici une recension
des symboles disponibles. Nous nous contenterons d'un avant goût des
principales catégories.

L'ouvrage de référence pour connaître les symboles disponibles dans
{\LaTeX} est la bien nommée \emph{Comprehensive {\LaTeX} Symbol List}
\citep{comprehensive}. La liste comprend près de \nombre{6000}
symboles répartis sur plus de 160 pages! On y trouve de tout, des
symboles mathématiques aux pictogrammes de météo ou d'échecs, en
passant par\dots\ des figurines des Simpsons.

\begin{important}
  Les moteurs {\XeTeX} et {\LuaTeX} supportent nativement le code
  source en format Unicode UTF-8 \citep{Unicode:5.0}. Ce standard
  contient des définitions pour plusieurs symboles mathématiques
  \citep{wikipedia:unicode-math}. Cela signifie qu'il est possible
  d'entrer une partie au moins des équations mathématiques avec des
  caractères visibles à l'écran, plutôt qu'avec des commandes
  {\LaTeX}. Nous ne saurions toutefois recommander cette pratique qui
  rend les fichiers source moins compatibles d'un système à un autre.
\end{important}

\subsection{Lettres grecques}
\label{sec:math:symboles:grecques}

On obtient les lettres grecques en {\LaTeX} avec des commandes
correspondant au nom de chaque lettre. Lorsque la commande débute par
une capitale, on obtient une lettre grecque majuscule. Les commandes
de certaines lettres grecques majuscules n'existent pas lorsque celles-ci
sont identiques aux lettres romaines.

Les tableaux \ref{tab:math:grecques} et \ref{tab:math:Grecques}
présentent l'ensemble des lettres grecques disponibles dans {\LaTeX}.

\begin{table}
  \caption{Lettres grecques minuscules}
  \label{tab:math:grecques}
  \begin{tabularx}{1.0\linewidth}{lXlXlXlX}
    $\alpha$      & \cmd{\alpha}      & $\theta$    & \cmd{\theta} &
    $o$           & o                 & $\tau$      & \cmd{\tau} \\
    $\beta$       & \cmd{\beta}       & $\vartheta$ & \cmd{\vartheta} &
    $\pi$         & \cmd{\pi}         & $\upsilon$  & \cmd{\upsilon} \\
    $\gamma$      & \cmd{\gamma}      & $\iota$     & \cmd{\iota} &
    $\varpi$      & \cmd{\varpi}      & $\phi$      & \cmd{\phi} \\
    $\delta$      & \cmd{\delta}      & $\kappa$    & \cmd{\kappa} &
    $\rho$        & \cmd{\rho}        & $\varphi$   & \cmd{\varphi} \\
    $\epsilon$    & \cmd{\epsilon}    & $\lambda$   & \cmd{\lambda} &
    $\varrho$     & \cmd{\varrho}     & $\chi$      & \cmd{\chi} \\
    $\varepsilon$ & \cmd{\varepsilon} & $\mu$       & \cmd{\mu} &
    $\sigma$      & \cmd{\sigma}      & $\psi$      & \cmd{\psi} \\
    $\zeta$       & \cmd{\zeta}       & $\nu$       & \cmd{\nu} &
    $\varsigma$   & \cmd{\varsigma}   & $\omega$    & \cmd{\omega} \\
    $\eta$        & \cmd{\eta}        & $\xi$       & \cmd{\xi}
  \end{tabularx}
\end{table}

\begin{table}
  \caption{Lettres grecques majuscules}
  \label{tab:math:Grecques}
  \begin{tabularx}{1.0\linewidth}{lXlXlXlX}
    $\Gamma$    & \cmd{\Gamma}   &
    $\Lambda$   & \cmd{\Lambda}  &
    $\Sigma$    & \cmd{\Sigma}   &
    $\Psi$      & \cmd{\Psi}     \\
    $\Delta$    & \cmd{\Delta}   &
    $\Xi$       & \cmd{\Xi}      &
    $\Upsilon$  & \cmd{\Upsilon} &
    $\Omega$    & \cmd{\Omega}   \\
    $\Theta$    & \cmd{\Theta}   &
    $\Pi$       & \cmd{\Pi}      &
    $\Phi$      & \cmd{\Phi}
  \end{tabularx}
\end{table}

\subsection{Lettres modifiées}
\label{sec:math:symboles:mathcal}

Les lettres de l'alphabet, principalement en majuscule, servent
parfois en mathématiques dans des versions modifiées pour représenter
des quantités, notamment les ensembles.

La commande \cmd{\mathcal} permet de transformer un ou plusieurs
caractères en version dite «calligraphique» dans le mode mathématique.
\begin{demo}
  \begin{texample}
\begin{lstlisting}
\mathcal{ABC\, xyz}
\end{lstlisting}
    \producing
    %% pris de lucidaot.tex; apparemment lent
    \setmathfont[RawFeature=+ss04]{Lucida Bright Math OT}
    $\mathscr{ABC\, xyz}$
  \end{texample}
\end{demo}

La commande \cmd{\mathbb} fournie par les paquetages \pkg{amsfonts} et
\pkg{unicode-math}, entre autres, génère des versions majuscule
ajourée (\emph{blackboard bold}) des lettres de l'alphabet. Elles sont
principalement utilisée pour représenter les ensembles de nombres.
\begin{demo}
  \begin{texample}
\begin{lstlisting}
\mathbb{NZRC}
\end{lstlisting}
    \producing $\mathbb{NZRC}$
  \end{texample}
\end{demo}

Le tableau~213 de la %
\doc[\emph{Comprehensive {\LaTeX} Symbol List}]{comprehensive}{http://texdoc.net/pkg/comprehensive} %
présente plusieurs autres alphabets spéciaux disponibles en mode
mathématique.

\begin{conseil}
  Certaines polices OpenType contiennent plusieurs
  versions des symboles mathématiques. Par exemple, la police utilisée
  dans le présent document contient deux versions de la police
  calligraphique, celle présentée ci-contre et celle-ci: %
  \setmathfont[RawFeature=-ss04]{Lucida Bright Math OT}
  $\mathscr{ABC\, xyz}$. Consulter éventuellement la documentation de
  la police pour les détails.
\end{conseil}

\subsection{Opérateurs binaires et relations}
\label{sec:math:symboles:binaires+relations}

Les opérateurs binaires combinent deux quantités pour en former une
troisième; pensons simplement aux opérateurs d'addition $+$ et de
soustraction $-$ que l'on retrouve sur un clavier d'ordinateur normal.
Les relations, quant à elles, servent pour la comparaison entre deux
quantités, comme $<$ et $>$. Le \autoref{tab:math:binaires}
présente une sélection d'opérateurs binaires et
le \autoref{tab:math:relations}  une sélection de relations.

La %
\doc[\emph{Comprehensive {\LaTeX} Symbol List}]{}{http://texdoc.net/pkg/comprehensive} %
consacre plus d'une dizaine de tableaux aux opérateurs binaires et
près d'une quarantaine aux relations. C'est dire à quel point les
tableaux \ref*{tab:math:binaires} et \ref*{tab:math:relations} de la
\autopageref{tab:math:binaires} ne présentent que les principaux
éléments à titre indicatif.

Certaines relations existent directement en version opposée, ou
négative (comme $\neq$ ou $\notin$) soit dans {\LaTeX} de base, soit
avec \pkg{amsmath} ou un autre paquetage. Autrement, il est possible
de préfixer toute relation de \cmd{\not} pour y superposer une barre
oblique $/$.

\begin{table}[p]
  \caption{Quelques opérateurs binaires}
  \label{tab:math:binaires}
  \begin{tabularx}{1.0\linewidth}{lXlXlXlX}
    $\times$    & \cmd{\times} &
    $\div$      & \cmd{\div}   &
    $\pm$       & \cmd{\pm}    &
    $\cdot$     & \cmd{\cdot}    \\
    $\cup$      & \cmd{\cup} &
    $\cap$      & \cmd{\cap} &
    $\setminus$ & \cmd{\setminus} &
    $\circ$     & \cmd{\circ}  \\
    $\wedge$    & \cmd{\wedge} &
    $\vee$      & \cmd{\vee} &
    $\oplus$    & \cmd{\oplus} &
    $\otimes$   & \cmd{\otimes} \\
    $\ast$      & \cmd{\ast} &
    $\star$     & \cmd{\star} &
    $\boxplus$  & \cmd{\boxplus}$^\dagger$ &
    $\boxtimes$ & \cmd{\boxtimes}$^\dagger$ \\
    \addlinespace
  \end{tabularx}
  \hspace*{1em}{\footnotesize $^\dagger$ requiert \pkg{amsmath}}
\end{table}

\begin{table}[p]
  \caption{Quelques relations et leur négation}
  \label{tab:math:relations}
  \begin{tabularx}{1.0\linewidth}{lXlXlXlX}
    $\leq$      & \cmd{\leq} &
    $\geq$      & \cmd{\geq}   &
    $\neq$      & \cmd{\neq}    &
    $\equiv$    & \cmd{\equiv}    \\
    $\subset$   & \cmd{\subset} &
    $\subseteq$ & \cmd{\subseteq}  &
    $\in$       & \cmd{\in} &
    $\notin$    & \cmd{\notin} \\
    $\nless$    & \cmd{\nless}$^\dagger$ &
    $\ngtr$     & \cmd{\ngtr}$^\dagger$   &
    $\nleq$     & \cmd{\nleq}$^\dagger$    &
    $\ngeq$     & \cmd{\ngeq}$^\dagger$ \\
    \addlinespace
  \end{tabularx}
  \hspace*{1em}{\footnotesize $^\dagger$ requiert \pkg{amsmath}}
\end{table}

\subsection{Flèches}
\label{sec:math:symboles:fleches}

Les flèches de différents types sont souvent utilisées en notation
mathématique, notamment dans les limites ou pour les expressions
logiques. Le \autoref{tab:math:fleches} en présente une sélection.

On retrouve les flèches utilisables en notation mathématique dans les
tableaux 102 à 119 de la %
\doc[\emph{Comprehensive {\LaTeX} Symbol List}]{}{http://texdoc.net/pkg/comprehensive}. %
Le document contient divers autres types de flèches, mais celles-ci ne
sont généralement pas appropriées pour les mathématiques (pensons à
{\manerrarrow} ou {\faArrowRight}).

Le paquetage \pkg{amsmath} fournit plusieurs flèches additionnelles
ainsi que la négation des plus communes. Ces dernières apparaissent
d'ailleurs dans le \autoref{tab:math:fleches}.

\begin{table}[p]
  \caption{Quelques flèches et leur négation}
  \label{tab:math:fleches}
  \begin{tabularx}{1.0\linewidth}{lXlX}
    $\gets$                & \cmd{\leftarrow}\quad \cmd{\gets} &
    $\longleftarrow$       & \cmd{\longleftarrow}              \\
    $\Leftarrow$           & \cmd{\Leftarrow}                  &
    $\Longleftarrow$       & \cmd{\Longleftarrow}              \\
    $\to$                  & \cmd{\rightarrow}\quad \cmd{\to}  &
    $\longrightarrow$      & \cmd{\longrightarrow}             \\
    $\Rightarrow$          & \cmd{\Rightarrow}                 &
    $\Longrightarrow$      & \cmd{\Longrightarrow}             \\
    $\uparrow$             & \cmd{\uparrow}                    &
    $\downarrow$           & \cmd{\downarrow}                  \\
    $\Uparrow$             & \cmd{\Uparrow}                    &
    $\Downarrow$           & \cmd{\Downarrow}                  \\
    $\updownarrow$         & \cmd{\updownarrow}                &
    $\Updownarrow$         & \cmd{\Updownarrow}                \\
    $\leftrightarrow$      & \cmd{\leftrightarrow}             &
    $\longleftrightarrow$  & \cmd{\longleftrightarrow}         \\
    $\Leftrightarrow$      & \cmd{\Leftrightarrow}             &
    $\Longleftrightarrow$  & \cmd{\Longleftrightarrow}         \\
    $\nleftarrow$          & \cmd{\nleftarrow}$^\dagger$         &
    $\nleftrightarrow$     & \cmd{\nleftrightarrow}$^\dagger$    \\
    $\nrightarrow$         & \cmd{\nrightarrow}$^\dagger$        &
    $\nLeftarrow$          & \cmd{\nLeftarrow}$^\dagger$         \\
    $\nLeftrightarrow$     & \cmd{\nLeftrightarrow}$^\dagger$    &
    $\nRightarrow$         & \cmd{\nRightarrow}$^\dagger$        \\
    \addlinespace
  \end{tabularx}
  \hspace*{1em}{\footnotesize $^\dagger$ requiert \pkg{amsmath}}
\end{table}

\subsection{Accents et autres symboles utiles}
\label{sec:math:symboles:autres}

Le \autoref{tab:math:autres} présente quelques uns des accents
disponibles dans le mode mathématique, ainsi que divers symboles
fréquemment utilisés en mathématiques.

Pour connaître l'ensemble des accents du mode mathématique de
{\LaTeX}, consulter le tableau~164 de la %
\doc[\emph{Comprehensive {\LaTeX} Symbol List}]{}{http://texdoc.net/pkg/comprehensive}. %
Les versions extensibles de certains accents se trouvent au
tableau~169. Quant aux symboles mathématiques divers, on en trouve de
toutes les sortes dans les tableaux 201--212.

\begin{table}[p]
  \caption{Accents et symboles mathématiques divers}
  \label{tab:math:autres}
  \begin{tabularx}{1.0\linewidth}{lXlXlXlX}
    $\hat{a}$ & \cmd{\hat}\verb={a}= &
    $\bar{a}$ & \cmd{\bar}\verb={a}= &
    $\tilde{a}$ & \cmd{\tilde}\verb={a}= &
    $\ddot{a}$ & \cmd{\ddot}\verb={a}= \\
    $\infty$ & \cmd{\infty} &
    $\nabla$ & \cmd{\nabla} &
    $\partial$ & \cmd{\partial} &
    $\ell$ & \cmd{\ell} \\
    $\forall$ & \cmd{\forall} &
    $\exists$ & \cmd{\exists} &
    $\emptyset$ & \cmd{\emptyset} &
    $\prime$ & \cmd{\prime} \\
    $\neg$ & \cmd{\neg} &
    $\backslash$ & \cmd{\backslash} &
    $\|$ & \pixbar &
    $\angle$ & \cmd{\angle}
  \end{tabularx}
\end{table}


\begin{exemple}
  L'équation suivante contient plusieurs des éléments présentés dans
  cette section et la précédente:
  \begin{equation*}
    \frac{\Gamma(\alpha)}{\lambda^\alpha} =
    \sum_{j = 0}^\infty \int_j^{j + 1} x^{\alpha - 1} e^{-\lambda x}\,
    dx,
    \quad
    \alpha > 0 \text{ et } \lambda > 0.
  \end{equation*}
  \begin{demo}
    \begin{texample}[0.58\linewidth]
\begin{lstlisting}
\begin{equation*}
\end{lstlisting}
      \producing
      \emph{équation hors paragraphe}
    \end{texample}

    \begin{texample}[0.58\linewidth]
\begin{lstlisting}
\frac{\Gamma(\alpha)}{
  \lambda^\alpha} =
\end{lstlisting}
      \producing
      $\dfrac{\Gamma(\alpha)}{\lambda^\alpha} =$
    \end{texample}

    \begin{texample}[0.58\linewidth]
\begin{lstlisting}
\sum_{j = 0}^\infty \int_j^{j + 1}
\end{lstlisting}
      \producing
      $\displaystyle \sum_{j = 0}^\infty \int_j^{j + 1}$
    \end{texample}

    \begin{texample}[0.58\linewidth]
\begin{lstlisting}
x^{\alpha - 1} e^{-\lambda x}\, dx
\end{lstlisting}
      \producing
      $\displaystyle x^{\alpha - 1} e^{-\lambda x}\, dx$
    \end{texample}

    \begin{texample}[0.58\linewidth]
\begin{lstlisting}
, \quad \alpha > 0 \text{ et }
\lambda > 0.
\end{lstlisting}
      \producing
      $, \quad \alpha > 0 \text{ et } \lambda > 0.$
    \end{texample}

    \begin{texample}[0.58\linewidth]
\begin{lstlisting}
\end{equation*}
\end{lstlisting}
      \producing
      \emph{fin de l'environnement}
    \end{texample}
  \end{demo}
  \qed
\end{exemple}


\section{Équations sur plusieurs lignes et numérotation}
\label{sec:math:align}

Dans ce qui précède, nous n'avons présenté que des équations tenant
sur une seule ligne en mode hors paragraphe. Cette section se penche
sur la manière de représenter des groupes d'équations du type
\begin{align}
  y &= 2x + 4 \\
  y &= 6x - 1
\end{align}
ou des suites d'équations comme
\begin{align*}
  x_{\text{max}}
  &= \sum_{i = 0}^{m - 1} (b - 1) b^i \\
  &= (b - 1) \sum_{i = 0}^{m - 1} b^i \\
  &= b^m - 1.
\end{align*}

Nous recommandons fortement les environnements de \pkg{amsmath} pour
les équations sur plusieurs lignes: ils sont plus polyvalents, plus
simples à utiliser et leur rendu est meilleur. Le
\autoref{tab:math:displays} --- repris presque intégralement de la
documentation de ce paquetage --- compare les différents
environnements pour les équations hors paragraphe.

\begin{table}[p]
  \caption{Comparaison des environnements pour les équations hors
    paragraphe de \pkg{amsmath} (les lignes verticales indiquent les
    marges logiques).}
  \label{tab:math:displays}
  \renewcommand{\theequation}{\arabic{equation}}
  \begin{eqxample}
\begin{lstlisting}
\begin{equation*}
  a = b
\end{equation*}
\end{lstlisting}
    \producing
    \begin{equation*}
      a = b
    \end{equation*}
  \end{eqxample}

  \begin{eqxample}
\begin{lstlisting}
\begin{equation}
  a = b
\end{equation}
\end{lstlisting}
    \producing
    \begin{equation}
      a = b
    \end{equation}
  \end{eqxample}

  \begin{eqxample}
\begin{lstlisting}
\begin{equation}
  \label{eq:5}
  \begin{split}
    a &= b + c - d \\
    &\phantom{=} + e - f \\
    &= g + h \\
    &= i
  \end{split}
\end{equation}
\end{lstlisting}
    \producing
    \begin{equation}\label{eq:math:5}
      \begin{split}
        a& =b+c-d\\
        &\phantom{=} +e-f\\
        & =g+h\\
        & =i
      \end{split}
    \end{equation}
  \end{eqxample}

  \begin{eqxample}
\begin{lstlisting}
\begin{multline}
  a + b + c + d + e + f \\
  + i + j + k + l + m + n
\end{multline}
\end{lstlisting}
    \producing
    \begin{multline}
      a+b+c+d+e+f\\
      +i+j+k+l+m+n
    \end{multline}
  \end{eqxample}

  \begin{eqxample}
\begin{lstlisting}
\begin{gather}
  a_1 = b_1 + c_1 \\
  a_2 = b_2 + c_2 - d_2 + e_2
\end{gather}
\end{lstlisting}
    \producing
    \begin{gather}
      a_1=b_1+c_1\\
      a_2=b_2+c_2-d_2+e_2
    \end{gather}
  \end{eqxample}

  \begin{eqxample}
\begin{lstlisting}
\begin{align}
  a_1 &= b_1 + c_1 \\
  a_2 &= b_2 + c_2 - d_2 + e_2
\end{align}
\end{lstlisting}
    \producing
    \begin{align}
      a_1& =b_1+c_1\\
      a_2& =b_2+c_2-d_2+e_2
    \end{align}
  \end{eqxample}

  \begin{eqxample}
\begin{lstlisting}
\begin{align}
  a_{11} &= b_{11} &
  a_{12} &= b_{12} \\
  a_{21} &= b_{21} &
  a_{22} &= b_{22} + c_{22}
\end{align}
\end{lstlisting}
    \producing
    \begin{align}
      a_{11}& =b_{11}&
      a_{12}& =b_{12}\\
      a_{21}& =b_{21}&
      a_{22}& =b_{22}+c_{22}
    \end{align}
  \end{eqxample}
\end{table}

\begin{itemize}
\item L'environnement de base pour les équations alignées sur un
  symbole de relation (en une ou plusieurs colonnes) est \Ie{align}.
  C'est l'environnement le plus utilisé en mode mathématique hormis
  \Pe{equation}.
\item Les environnements \Pe{multline}, \Pe{gather} et \Pe{align}
  existent également en version étoilée (\Ie{multline*}, \Ie{gather*},
  \Ie{align*}) qui ne numérotent pas les équations.
\item Dans les environnements \Pe{align} et \Ie{split}, les équations
  successives sont alignées sur le caractère se trouvant immédiatement
  après le marqueur de colonne \verb=&=.
\item Comme dans les tableaux, la commande {\pixbsbs} sert à
  délimiter les lignes de la suite d'équations.
\item Remarquer, dans le troisième exemple du
  \autoref{tab:math:displays}, comment la commande \cmd{\phantom} sert
  à insérer un blanc exactement de la largeur du symbole $=$ au début
  de la seconde ligne de la suite d'égalités.
\item Pour supprimer la numérotation d'une ligne dans une série
  d'équations numérotées, placer la commande \cmd{\notag} juste avant
  la commande de changement de ligne {\bs\bs}.
  \begin{demo}
    \begin{texample}[0.53\linewidth]
\begin{lstlisting}
\begin{align}
  a_1 &= b_1 + c_1 \notag \\
  a_2 &= b_2 + c_2 - d_2 + e_2
\end{align}
\end{lstlisting}
      \producing
      \begin{align}
        a_1& =b_1+c_1 \notag \\
        a_2& =b_2+c_2-d_2+e_2
      \end{align}
    \end{texample}
  \end{demo}
\item Les renvois vers des équations numérotées fonctionnent, comme
  partout ailleurs en {\LaTeX}, avec le système d'étiquettes et de
  références (\autoref{sec:tableaux:floats}). Le paquetage
  \pkg{amsmath} fournit également la pratique commande \cmd{\eqref}
  qui place automatiquement le numéro d'équation entre parenthèses.
  \begin{demo}
    \begin{texample}
\begin{lstlisting}
L'équation \eqref{eq:5} du
tableau `\ref*{tab:math:displays}' démontre que...
\end{lstlisting}
      \producing
      L'équation \eqref{eq:math:5} du tableau
      \ref*{tab:math:displays} démontre que\dots
    \end{texample}
  \end{demo}
\item L'environnement \Ie{split} sert à apposer un seul numéro à une
  équation affichée sur plusieurs lignes. Il doit être employé à
  l'intérieur d'un autre environnement d'équations hors paragraphe.
\end{itemize}
Consulter le chapitre 3 de la %
\doc{amsldoc}{http://texdoc.net/pkg/amsmath} %
du paquetage \pkg{amsmath} pour les détails sur
l'utilisation des environnements du \autoref*{tab:math:displays}.

\begin{conseil}
  Veillez à respecter les règles suivantes pour la composition des
  équations.
  \begin{enumerate}
  \item Qu'elles apparaissent en ligne ou hors paragraphe, les
    équations font partie intégrante de la phrase. Ainsi, les règles
    de ponctuation usuelles s'appliquent-elles aux équations.
  \item Lorsqu'une équation s'étend sur plus d'une ligne, couper
    chaque ligne \emph{avant} un opérateur de sorte que chaque ligne
    constitue une expression mathématique complète (voir les troisième
    et quatrième exemples du \autoref{tab:math:displays}).
  \item Ne numéroter que les équations d'un document auxquelles le
    texte fait référence.
  \end{enumerate}
\end{conseil}

\begin{exemple}
  Nous avons réalisé les deux suites d'équations au début de la
  section avec les extraits de code ci-dessous, dans l'ordre.
\begin{lstlisting}
\begin{align}
  y &= 2x + 4 \\
  y &= 6x - 1
\end{align}
\end{lstlisting}
\begin{lstlisting}
\begin{align*}
  x_{\text{max}}
  &= \sum_{i = 0}^{m - 1} (b - 1) b^i \\
  &= (b - 1) \sum_{i = 0}^{m - 1} b^i \\
  &= b^m - 1.
\end{align*}
\end{lstlisting}
  \qed
\end{exemple}


\section{Délimiteurs de taille variable}
\label{sec:math:delimiteurs}

Les délimiteurs en mathématiques sont des symboles généralement
utilisés en paire tels que les parenthèses $(~)$, les crochets $[~]$
ou les accolades $\{~\}$ et qui servent à regrouper des termes d'une
équation. La taille des délimiteurs doit s'adapter au contenu entre
ceux-ci afin d'obtenir, par exemple,
\begin{equation*}
  \left( 1 + \frac{1}{x} \right)
\end{equation*}
plutôt que la peu esthétique composition
\begin{equation*}
  (1 + \frac{1}{x}).
\end{equation*}

La paire de commandes
\begin{lstlisting}
\left`\meta{delim\_g}'  ...  \right`\meta{delim\_d}'
\end{lstlisting}
définit un délimiteur gauche \meta{delim\_g} et un délimiteur droit
\meta{delim\_d} dont la taille s'ajustera automatiquement au contenu
entre les deux commandes.
\begin{demo}
  \begin{texample}
\begin{lstlisting}
\left( 1 + \frac{1}{x} \right)
\end{lstlisting}
    \producing
    \begin{equation*}
      \left( 1 + \frac{1}{x} \right)
    \end{equation*}
  \end{texample}

  \begin{texample}
\begin{lstlisting}
\left(
  \sum_{i = 1}^n x_i^2
\right)^{1/2}
\end{lstlisting}
    \producing
    \begin{equation*}
      \left(
        \sum_{i = 1}^n x_i^2
      \right)^{1/2}
    \end{equation*}
  \end{texample}
\end{demo}
Les commandes \cmd{\left} et \cmd{\right} doivent toujours former une
paire, c'est-à-dire qu'à \emph{toute} commande \cmdprint{\left} doit
absolument correspondre une commande \cmdprint{\right}. Cette
contrainte est facile à oublier!

\begin{itemize}
\item Il est possible d'imbriquer des paires de commandes les unes à
  l'intérieur des autres, pour autant que l'expression compte toujours
  autant de \cmdprint{\left} que de \cmdprint{\right}.
  \begin{demo}
    \begin{texample}
\begin{lstlisting}
\left[
  \int
  \left(
    1 + \frac{x}{k}
  \right) dx
\right]
\end{lstlisting}
      \producing
      \begin{equation*}
        \left[
          \int
          \left(
            1 + \frac{x}{k}
          \right) dx
        \right]
      \end{equation*}
    \end{texample}
  \end{demo}
%
\item Les symboles \meta{delim\_g} et \meta{delim\_d} n'ont pas à
  former une paire logique; toute combinaison de délimiteurs est
  valide.
  \begin{demo}
    \begin{texample}
\begin{lstlisting}
\int_0^1 x\, dx =
\left[
  \frac{x^2}{2}
\right|_0^1
\end{lstlisting}
      \producing
      \begin{equation*}
        \int_0^1 x\, dx =
        \left[
          \frac{x^2}{2}
        \right|_0^1
      \end{equation*}
    \end{texample}
  \end{demo}
%
\item Il arrive qu'un seul délimiteur soit nécessaire. Pour respecter
  la règle de la paire ci-dessus, on aura recours dans ce cas à un
  délimiteur \emph{invisible} représenté par le caractère «\verb=.=».
  \begin{demo}
    \begin{texample}
\begin{lstlisting}
f(x) =
\left\{
  \begin{aligned}
    1 - x, &\quad x < 1 \\
    x - 1, &\quad x \geq 1
  \end{aligned}
\right.
\end{lstlisting}
      \producing
      \begin{equation*}
        f(x) =
        \left\{
          \begin{aligned}
            1 - x, &\quad x < 1 \\
            x - 1, &\quad x \geq 1
          \end{aligned}
        \right.
      \end{equation*}
    \end{texample}
  \end{demo}
  (L'environnement \Ie{aligned} utilisé ci-dessus provient de
  \pkg{amsmath}.) On notera au passage que l'environnement \Ie{cases}
  de \pkg{amsmath} rend plus simple la réalisation de constructions
  comme celle ci-dessus.
  \begin{demo}
    \begin{texample}
\begin{lstlisting}
f(x) =
\begin{cases}
  1 - x, & x < 1 \\
  x - 1, & x \geq 1
\end{cases}
\end{lstlisting}
      \producing
      \begin{equation*}
        f(x) =
        \begin{cases}
          1 - x, & x < 1 \\
          x - 1, & x \geq 1
        \end{cases}
      \end{equation*}
    \end{texample}
  \end{demo}
%
\item La règle de la paire est tout spécialement délicate dans les
  équations sur plusieurs lignes car elle s'applique à chaque ligne
  d'une équation. Par conséquent, si la paire de délimiteurs s'ouvre
  sur une ligne et se referme sur une autre, il faudra ajouter un
  délimiteur invisible à la fin de la première ligne ainsi qu'au début
  de la seconde.
  \begin{demo}
    \begin{texample}
\begin{lstlisting}
\begin{align*}
  a
  &= \left(
    b + \frac{c}{d}
    \right. \\
  &\phantom{=} \left.
    + \frac{e}{d} - f
    \right)
\end{align*}
\end{lstlisting}
      \producing
      \begin{align*}
        a
        &= \left(
          b + \frac{c}{d}
          \right. \\
        &\phantom{=} \left.
          + \frac{e}{d} - f
          \right)
      \end{align*}
    \end{texample}
  \end{demo}
%
\item Quand les choix de taille de délimiteurs de {\LaTeX} ne
  conviennent pas pour une raison ou pour une autre, on peut
  sélectionner leur taille avec les commandes %
  \cmd{\big}, %
  \cmd{\Big}, %
  \cmd{\bigg} et %
  \cmd{\Bigg}. %
  Ces commandes s'utilisent comme \cmdprint{\left} et
    \cmdprint{\right} en les faisant immédiatement suivre d'un
  délimiteur. Le \autoref{tab:math:big_et_al} contient des exemples de
  délimiteurs pour chaque taille.

  \begin{table}
    \centering
    \caption{Tailles des délimiteurs mathématiques}
    \label{tab:math:big_et_al}
    \begin{tabular}{ll}
      \toprule
      taille standard & $(~) \quad [~] \quad \{~\}$ \\
      \addlinespace[0.5\normalbaselineskip]
      \cmd{\big} & $\big(~\big) \quad \big[~\big] \quad \big\{~\big\}$ \\
      \addlinespace[0.5\normalbaselineskip]
      \cmd{\Big} & $\Big(~\Big) \quad \Big[~\Big] \quad \Big\{~\Big\}$ \\
      \addlinespace[0.5\normalbaselineskip]
      \cmd{\bigg} & $\bigg(~\bigg) \quad \bigg[~\bigg] \quad \bigg\{~\bigg\}$ \\
      \addlinespace[0.5\normalbaselineskip]
      \cmd{\Bigg} & $\Bigg(~\Bigg) \quad \Bigg[~\Bigg] \quad \Bigg\{~\Bigg\}$ \\
      \addlinespace[4pt]\bottomrule
    \end{tabular}
  \end{table}
\end{itemize}

La section~14 de la %
\doc{amsldoc}{http://texdoc.net/pkg/amsmath} %
de \pkg{amsmath} traite de divers enjeux typographiques en lien
avec les délimiteurs et on y introduit des nouvelles commandes pour
contrôler leur taille. C'est une lecture suggérée.

\begin{exemple}
  Le développement de la formule d'approximation de Simpson comporte
  plusieurs des éléments discutés jusqu'à maintenant:
  \begin{align*}
    \int_a^b f(x)\, dx
    &\approx \sum_{j = 0}^{n - 1}
      \int_{x_{2j}}^{x_{2(j + 1)}} f(x)\, dx \\
    &= \frac{h}{3} \sum_{j = 0}^{n - 1}
      \left[
      f(x_{2j}) + 4 f(x_{2j + 1}) + f(x_{2(j + 1)})
      \right]
    \displaybreak[0] \\
    &= \frac{h}{3}
      \left[
      f(x_0) +
      \sum_{j = 1}^{n - 1} f(x_{2j}) +
      4 \sum_{j = 0}^{n - 1} f(x_{2j + 1})
      \right. \displaybreak[0] \\
    &\phantom{=}  + \left.
      \sum_{j = 0}^{n - 2} f(x_{2(j + 1)}) +
      f(x_{2n})
      \right] \\
    &= \frac{h}{3}
      \left[
      f(a) +
      2 \sum_{j = 1}^{n - 1} f(x_{2j}) +
      4 \sum_{j = 0}^{n - 1} f(x_{2j + 1}) +
      f(b)
      \right].
  \end{align*}
  On compose ce bloc d'équations avec le code source ci-dessous.
\begin{lstlisting}
\begin{align*}
  \int_a^b f(x)\, dx
  &\approx \sum_{j = 0}^{n - 1}
    \int_{x_{2j}}^{x_{2(j + 1)}} f(x)\, dx \\
  &= \frac{h}{3} \sum_{j = 0}^{n - 1}
    \left[
    f(x_{2j}) + 4 f(x_{2j + 1}) + f(x_{2(j + 1)})
    \right] \\
  &= \frac{h}{3}
    \left[
    f(x_0) +
    \sum_{j = 1}^{n - 1} f(x_{2j}) +
    4 \sum_{j = 0}^{n - 1} f(x_{2j + 1})
    \right. \\
  &\phantom{=} + \left.
    \sum_{j = 0}^{n - 2} f(x_{2(j + 1)}) +
    f(x_{2n})
    \right] \\
  &= \frac{h}{3}
    \left[
    f(a) +
    2 \sum_{j = 1}^{n - 1} f(x_{2j}) +
    4 \sum_{j = 0}^{n - 1} f(x_{2j + 1}) +
    f(b)
    \right].
\end{align*}
\end{lstlisting}
  \qed
\end{exemple}

\begin{conseil}
  De manière générale, il est déconseillé de scinder une suite
  d'équations entre deux pages. Le chargement du paquetage
  \pkg{amsmath} désactive d'ailleurs cette possibilité. Cependant,
  c'est parfois inévitable pour les longs blocs d'équations.

  La commande \cmd{\displaybreak}, placée immédiatement avant {\bs\bs}
  dans un environnement comme \Ie{align} permet d'indiquer à {\LaTeX}
  la possibilité d'insérer un saut de page après la ligne courante
  dans le bloc d'équations.

  La commande accepte en argument optionnel un entier entre $0$ et $4$
  indiquant à quel point un saut de page est désiré:
  \cmdprint{\displaybreak[0]} signifie «il est permis de changer de
  page ici» sans que ce ne soit obligatoire;
  \cmdprint{\displaybreak[4]} ou, de manière équivalente,
  \cmdprint{\displaybreak} force un saut de page.

  À moins d'en être vraiment aux toutes dernières étapes d'édition
  d'un document, il est conseillé d'utiliser la commande
  \cmdprint{\displaybreak} avec parcimonie et avec un argument
  optionnel faible.
\end{conseil}


\section{Caractères gras en mathématiques}
\label{sec:math:gras}

Les caractères gras sont parfois utilisés en mathématiques,
particulièrement pour représenter les vecteurs et les matrices:
\begin{equation*}
  \symbf{A} \symbf{x} = \symbf{b} \Leftrightarrow
  \symbf{x} = \symbf{A}^{-1} \symbf{b}
\end{equation*}

Pourquoi consacrer une section spécialement à cette convention
typographique? Parce que la création de symboles mathématiques en
gras doit certainement figurer parmi les questions les plus
fréquemment posées par les utilisateurs de {\LaTeX}\dots\ et que la
réponse n'est pas unique!

La commande
\begin{lstlisting}
\mathbf`\marg{symbole}'
\end{lstlisting}
place \meta{symbole} en caractère gras en mode mathématique. C'est
donc l'analogue de la commande \cmd{\textbf} du mode texte. Dans
{\LaTeX} de base, la commande n'a toutefois un effet que sur les
lettres latines et, parfois, les lettres grecques majuscules.
\begin{demo}
  \begin{texample}[0.6\linewidth]
\begin{lstlisting}
\theta \mathbf{\theta} +
\Gamma \mathbf{\Gamma} \mathbf{+}
A \mathbf{A}
\end{lstlisting}
    \producing
    %% ok, ici il faut vraiment tricher pour reproduire avec
    %% unicode-math ce qui se produirait avec pdflatex standard
    $\theta {\theta} + \Gamma \symbf{\Gamma} {+} A \symbfup{A}$
  \end{texample}
\end{demo}
On remarquera aussi que \verb=\mathbf{A}= produit une lettre
majuscule droite plutôt qu'en italique.

La manière la plus standard et robuste d'obtenir des symboles
mathématiques (autres que les lettres) en gras semble être, au moment
d'écrire ces lignes, via la commande
\begin{lstlisting}
\bm`\marg{symbole}'
\end{lstlisting}
fournie par le paquetage \pkg{bm} \citep{bm}.
\begin{demo}
  \begin{texample}[0.6\linewidth]
\begin{lstlisting}
\theta \bm{\theta} +
\Gamma \bm{\Gamma} \bm{+}
A \bm{A}
\end{lstlisting}
    \producing
    %% ok, ici il faut vraiment tricher pour reproduire avec
    %% unicode-math ce qui se produirait avec pdflatex standard
    $\theta \symbf{\theta} + \Gamma \symbfup{\Gamma} \symbf{+} A
    \symbf{A}$
  \end{texample}
\end{demo}

Les utilisateurs de {\XeLaTeX} chargent normalement le paquetage
\pkg{unicode-math} \citep{unicode-math} pour sélectionner la police
pour les mathématiques. Ce paquetage fournit la commande
\begin{lstlisting}
\symbf`\marg{symbole}'
\end{lstlisting}
pour placer un \meta{symbole} mathématique en gras. Le paquetage offre
différentes combinaisons de lettres latines et grecques droites ou
italiques en gras selon la valeur de l'option \code{bold-style};
consulter la section~5 de la %
\doc{unicode-math}{http://texdoc.net/pkg/unicode-math}. %
\begin{demo}
  \begin{texample}[0.6\linewidth]
\begin{lstlisting}
% XeLaTeX + paquetage unicode-math
% avec l'option bold-style=ISO
\theta \symbf{\theta} +
\Gamma \symbf{\Gamma} \symbf{+}
A \symbf{A}
\end{lstlisting}
    \producing
    $\theta \symbf{\theta} + \Gamma \symbf{\Gamma} \symbf{+} A
    \symbf{A}$
  \end{texample}
\end{demo}

\begin{conseil}
  \label{conseil:math:mat}
  Si le gras est fréquemment utilisé dans un document pour une
  notation particulière, il est fortement recommandé de définir une
  nouvelle commande\footnotemark\
  sémantique plutôt que d'utiliser à répétition l'une ou l'autre des
  commandes ci-dessus.

  Par exemple, si le gras est utilisé pour les vecteurs et matrices,
  on pourrait définir une nouvelle commande \cmdprint{\mat} en
  insérant dans le préambule du document
\begin{lstlisting}
\newcommand[1]{\mat}{\bm{#1}}
\end{lstlisting}
\end{conseil}
%
\footnotetext{%
  La définition de nouvelles commandes est couvert plus en détail au
  \autoref{chap:commandes}.}


\begin{information}
  On l'a vu ci-dessus: le paquetage \pkg{unicode-math} offre l'option
  \code{ISO} pour le traitement du gras dans les mathématiques. En
  fait, c'est toute la composition des mathématiques qui fait l'objet
  d'un standard ISO!

  Plusieurs prescriptions du standard --- comme les noms de variables
  en italique --- sont déjà prises en compte par {\LaTeX}. Le respect
  de certaines autres règles est moins répandu, notamment celle qui
  veut que les constantes mathématiques dont la valeur de change pas
  (ce sont des constantes, après tout) soient représentées en police
  droite. On pense, par exemple, au nombre d'Euler $\mathrm{e}$, au
  nombre imaginaire $\mathrm{i} = \sqrt{-1}$. Même chose pour les
  opérateurs mathématiques, dont le $\mathrm{d}$ dans les dérivées et
  les intégrales.

  Voici un exemple d'intégrale composée avec ces règles:
  \begin{equation*}
    \int_C \frac{\mathrm{e}^z}{z}\, \mathrm{d}z = 2 \pi \mathrm{i}.
  \end{equation*}
  Pour en savoir plus sur l'utilisation du standard ISO dans {\LaTeX},
  consulter \citet{Beccari:ISO} ou, pour un sommaire rapide, l'%
  \link{https://nickhigham.wordpress.com/2016/01/28/typesetting-mathematics-according-to-the-iso-standard/}{entrée de blogue} %
  de Nick Higham.
\end{information}

\begin{exemple}
  \label{ex:math:matrices}
  Le paquetage \pkg{amsmath} fournit quelques environnements qui
  facilitent la mise en forme de matrices; ils diffèrent simplement
  par le type de délimiteur autour de la matrice.

  Considérer le bloc d'équations suivant:
  \begin{align*}
    \mat{J}(\mat{\theta})
    &=
      \begin{bmatrix}
        \dfrac{\partial f_1(\mat{\theta})}{\partial \theta_1} &
        \dfrac{\partial f_1(\mat{\theta})}{\partial \theta_2} \\[12pt]
        \dfrac{\partial f_2(\mat{\theta})}{\partial \theta_1} &
        \dfrac{\partial f_2(\mat{\theta})}{\partial \theta_2}
      \end{bmatrix} \\
    &=
      \left[
      \frac{\partial f_i(\mat{\theta})}{\partial \theta_j}
      \right]_{2 \times 2}, \quad i, j = 1, 2.
  \end{align*}
  Le code pour composer ces équations est:
\begin{lstlisting}
\begin{align*}
  \mat{J}(\mat{\theta})
  &=
  \begin{bmatrix}
    \dfrac{\partial f_1(\mat{\theta})}{\partial \theta_1} &
    \dfrac{\partial f_1(\mat{\theta})}{\partial \theta_2}
    \\[12pt] % augmenter l'espace entre les lignes
    \dfrac{\partial f_2(\mat{\theta})}{\partial \theta_1} &
    \dfrac{\partial f_2(\mat{\theta})}{\partial \theta_2}
  \end{bmatrix} \\
  &=
  \left[
  \frac{\partial f_i(\mat{\theta})}{\partial \theta_j}
  \right]_{2 \times 2}, \quad i, j = 1, 2.
\end{align*}
\end{lstlisting}
  On remarquera l'utilisation de la commande \cmd{\dfrac}
  (\autoref{sec:math:bases:fractions}) pour composer des grandes
  fractions à l'intérieur des matrices. Nous avons également eu
  recours à la commande \cmdprint{\mat} mentionnée dans la
  rubrique Conseil du {\TeX}pert de la \autopageref{conseil:math:mat}. %
  \qed
\end{exemple}


%%%
%%% Exercices
%%%

\section{Exercices}
\label{sec:math:exercices}

\Opensolutionfile{solutions}[solutions-mathematiques]

\begin{Filesave}{solutions}
\section*{Chapitre \ref*{chap:math}}
\addcontentsline{toc}{section}{Chapitre \protect\ref*{chap:math}}

\end{Filesave}

\begin{exercice}
  Utiliser le gabarit de document \fichier{exercice\_gabarit.tex} pour
  reproduire le texte suivant:
  \begin{quote}
    La dérivée de la fonction composée $f \circ g(x) = f[g(x)]$ est
    $\{f[g(x)]\}^\prime = f^\prime[g(x)] g^\prime(x)$. La dérivée
    seconde du produit des fonctions $f$ et $g$ est
    \begin{equation*}
      [f(x) g(x)]^{\prime\prime} =
      f^{\prime\prime}(x) g(x) + 2 f^\prime(x) g^\prime(x) +
      f(x) g^{\prime\prime}(x).
    \end{equation*}
  \end{quote}
  \begin{sol}
    On trouve la commande pour produire le symbole $\circ$ dans le
    \autoref{tab:math:binaires}. On peut produire les symboles de
    dérivée $\prime$ avec la commande \cmd{\prime}
    (\autoref{tab:math:autres}) ou simplement avec le caractère
    \verb='=.
\begin{lstlisting}
La dérivée de la fonction composée $f \circ g(x) =
f[g(x)]$ est $\{f[g(x)]\}^\prime = f^\prime[g(x)]
g^\prime(x)$. La dérivée seconde du produit des
fonctions $f$ et $g$ est
\begin{equation*}
  [f(x) g(x)]^{\prime\prime} =
  f^{\prime\prime}(x) g(x) + 2 f^\prime(x) g^\prime(x) +
  f(x) g^{\prime\prime}(x).
\end{equation*}
\end{lstlisting}
  \end{sol}
\end{exercice}

\begin{exercice}
  Composer l'équation suivante avec l'environnement \Pe{align*}:
  \begin{align*}
    f(x +& h, y + k)
    = f(x, y) +
      \left\{
      \frac{\partial f(x, y)}{\partial x} h +
      \frac{\partial f(x, y)}{\partial y} k
      \right\} \\
    &+
      \frac{1}{2}
      \left\{
      \frac{\partial^2 f(x, y)}{\partial x^2} h^2 +
      \frac{\partial^2 f(x, y)}{\partial x \partial y} kh +
      \frac{\partial^2 f(x, y)}{\partial y^2} k^2
      \right\} \\
    &+
      \frac{1}{6} \{\cdots\} + \dots + \frac{1}{n!} \{\cdots\} + R_n.
  \end{align*}
  Aligner les deuxième et troisième lignes de l'équation sur divers
  caractères de la première ligne afin que l'équation ne dépasse pas
  les marges du document.
  \begin{sol}
    Nous avons aligné les lignes de l'équation juste à droite du
    premier symbole $+$ à la première ligne. Remarquer l'usage des
    commandes \cmd{\cdots} et \cmd{\dots} dans la dernière ligne:
    {\LaTeX} choisit correctement la position centrée des points entre les
    opérateurs d'addition, mais pas entre les accolades.
\begin{lstlisting}
\begin{align*}
  f(x +& h, y + k) = f(x, y) +
    \left\{
    \frac{\partial f(x, y)}{\partial x} h +
    \frac{\partial f(x, y)}{\partial y} k
    \right\} \\
  &+
    \frac{1}{2}
    \left\{
    \frac{\partial^2 f(x, y)}{\partial x^2} h^2 +
    \frac{\partial^2 f(x, y)}{\partial x \partial y} kh +
    \frac{\partial^2 f(x, y)}{\partial y^2} k^2
    \right\} \\
  &+
    \frac{1}{6} \{\cdots\} + \dots +
    \frac{1}{n!} \{\cdots\} + R_n.
\end{align*}
\end{lstlisting}
  \end{sol}
\end{exercice}

\begin{exercice}
  Composer à l'aide de l'environnement \Ie{cases}
  (\autoref{sec:math:delimiteurs}) la définition de la fonction
  $\tilde{f}(x)$:
  \begin{equation*}
    \tilde{f}(x) =
    \begin{cases}
      0, & x \leq c_0\\
      \dfrac{F_n(c_j) - F_n(c_{j-1})}{c_j - c_{j-1}} =
      \dfrac{n_j}{n (c_j - c_{j - 1})}, &  c_{j-1} < x \leq c_j\\
      0, & x > c_r.
    \end{cases}
  \end{equation*}
  Il est nécessaire d'imposer la taille des fractions dans la seconde
  branche de la définition à l'aide des fonctions de la
  \autoref{sec:math:bases:fractions}.
  \begin{sol}
    Il faut utiliser \cmd{\dfrac} pour obtenir des fractions dans une
    branche de \Pe{cases} de la même taille que dans une équation hors
    paragraphe:
\begin{lstlisting}
\begin{equation*}
  \tilde{f}(x) =
  \begin{cases}
    0, & x \leq c_0\\
    \dfrac{F_n(c_j) - F_n(c_{j-1})}{c_j - c_{j-1}} =
    \dfrac{n_j}{n (c_j - c_{j - 1})}, &
      c_{j-1} < x \leq c_j \\
    0, & x > c_r.
  \end{cases}
\end{equation*}
\end{lstlisting}
  \end{sol}
\end{exercice}

\begin{exercice}[nosol]
  Le fichier \fichier{exercice\_mathematiques.tex} contient un exemple
  complet de développement mathématique. Étudier le contenu du fichier
  puis compiler celui-ci tel quel avec {pdf\LaTeX} ou {\XeLaTeX}.
  Effectuer ensuite les modifications suivantes.
  \begin{enumerate}
  \item Charger le paquetage \pkg{amsfonts} dans le préambule, puis
    remplacer \verb=$R^+$= par \verb=$\mathbb{R}^+$= à la ligne
    débutant par «Le domaine».
  \item Dans l'équation du Jacobien de la transformation, remplacer
    successivement l'environnement \Ie{vmatrix} par %
    \Ie{pmatrix}, %
    \Ie{bmatrix}, %
    \Ie{Bmatrix} et %
    \Ie{Vmatrix}. Observer l'effet sur les délimiteurs de la matrice.
  \item Toujours dans la même matrice, composer successivement les
    deux fractions avec les commandes \cmd{\frac}, \cmd{\tfrac} et
    \cmd{\dfrac}. Observer le résultat.
  \item Réduire l'espacement de part et d'autre du symbole
    $\Leftrightarrow$ dans la seconde équation hors paragraphe.
  \item À l'aide de la fonction Rechercher et remplacer de l'éditeur
    de texte, remplacer toutes les occurrences du symbole $\theta$ par
    $\lambda$.
  \end{enumerate}
\end{exercice}

\Closesolutionfile{solutions}

%%% Local Variables:
%%% mode: latex
%%% TeX-engine: xetex
%%% TeX-master: "formation-latex-ul"
%%% coding: utf-8
%%% End:

version https://git-lfs.github.com/spec/v1
oid sha256:bde34e0d67701852574dce9d46b6551edf73017845c601bdb936b6e993de1b93
size 34993

version https://git-lfs.github.com/spec/v1
oid sha256:3cab9dad9c9aa96de811e3293fec99d28d04f509f2d40a058a663b9c49a40560
size 17069

\chapter{Trucs et astuces divers}
\label{chap:trucs}

En clôture de l'ouvrage, ce chapitre traite de différents sujets que
même une personne débutant avec {\LaTeX} voudra assez rapidement
aborder, comme le contrôle des sauts de ligne et des sauts de page, la
modification de la police du document, l'utilisation de la couleur ou
l'insertion d'hyperliens dans le fichier de sortie PDF. Nous offrons
également de courtes introductions à des usages plus spécialisés de
{\LaTeX} comme la mise en page de code informatique, la production de
diapositives ou la programmation lettrée. Enfin, nous expliquons
sommairement comme assurer de manière efficace la gestion des versions
de ses documents, surtout dans un contexte de travail collaboratif.

\section{Contrôle de la disposition du texte}
\label{sec:trucs:controle}

Des algorithmes élaborés permettent généralement à {\LaTeX} de
disposer le texte harmonieusement sur la page. Néanmoins, il
est parfois nécessaire  d'effectuer soi-même de menus ajustements,
notamment pour les sauts de page et la coupure de mots.

\subsection{Sauts de ligne et de page}
\label{sec:trucs:controle:sauts}

Il est assez rarement nécessaire avec {\LaTeX} de devoir forcer les
retours à la ligne. Chose certaine, l'on devrait toujours utiliser une
ligne blanche dans le code source pour identifier un changement de
paragraphe.

Cela dit, les commandes suivantes permettent d'insérer un saut de
ligne manuellement lorsque requis:
\begin{lstlisting}
\\
\\`\oarg{longueur}'
\newline
\end{lstlisting}
La commande {\pixbsbs} est connue: elle sert aussi à délimiter les
lignes dans les tableaux (\autoref{sec:tableaux:tableaux}) et les lignes
d'une suite d'équations (\autoref{sec:math:align}). L'argument optionnel
\oarg{longueur} permet d'insérer un blanc entre les deux lignes; la
\autoref{sec:bases:longueurs} explique comment spécifier une
longueur.

Généralement équivalente à {\bs\bs}, la commande \cmd{\newline} est
parfois nécessaire, notamment pour insérer un changement de ligne à
l'intérieur d'une cellule d'un tableau ou à l'intérieur d'un titre de
section. Quand {\bs\bs} ne fonctionne pas, essayer
\cmdprint{\newline}.

\begin{exemple}
  La commande {\bs\bs} est particulièrement utile --- voire nécessaire
  --- pour disposer des boîtes à l'intérieur d'une figure.
  L'utilisation de l'argument \meta{longueur} permet alors de
  contrôler l'espacement vertical entre les éléments. Comparer les
  deux exemples ci-dessous.
  \begin{demo}
    \begin{texample}[0.55\linewidth]
\begin{lstlisting}
\begin{minipage}{1.0\linewidth}
  \framebox[\linewidth]{texte}
\end{minipage} \\
\begin{minipage}{1.0\linewidth}
  \framebox[\linewidth]{texte}
\end{minipage}
\end{lstlisting}
      \producing
      \begin{minipage}{1.0\linewidth}
        \framebox[\linewidth]{texte}
      \end{minipage}
      \begin{minipage}{1.0\linewidth}
        \framebox[\linewidth]{texte}
      \end{minipage}
    \end{texample}

    \begin{texample}[0.55\linewidth]
\begin{lstlisting}
\begin{minipage}{1.0\linewidth}
  \framebox[\linewidth]{texte}
\end{minipage} \\`\textbf{[6pt]}'
\begin{minipage}{1.0\linewidth}
  \framebox[\linewidth]{texte}
\end{minipage}
\end{lstlisting}
      \producing
      \begin{minipage}{1.0\linewidth}
        \framebox[\linewidth]{texte}
      \end{minipage} \\[6pt]
      \begin{minipage}{1.0\linewidth}
        \framebox[\linewidth]{texte}
      \end{minipage}
    \end{texample}
  \end{demo}
  \qed
\end{exemple}

Les commandes
\begin{lstlisting}
\newpage
\clearpage
\cleartorecto              % classe memoir seulement
\cleartoverso              % classe memoir seulement
\end{lstlisting}
permettent d'insérer manuellement un saut de page pour éviter une
coupure malheureuse. La commande de base pour insérer un saut
n'importe où dans la page est \cmd{\newpage}. La commande
\cmd{\clearpage}, quant à elle, va également s'assurer d'afficher tous
les éléments flottants (\autoref{sec:tableaux:floats}) en attente de
disposition.

Les commandes \cmd{\cleartorecto} et \cmd{\cleartoverso}, propres à la
classe \class{memoir} permettent respectivement de passer
automatiquement à une page recto ou à une page verso. Évidemment, elles
n'ont d'utilité que dans les documents recto verso.

Moins directives, les commandes
\begin{lstlisting}
\pagebreak`\oarg{n}'
\enlargethispage`\marg{longueur}'
\end{lstlisting}
permettent de seulement aider {\LaTeX} à gérer les sauts de page à un
endroit précis. La commande \cmd{\pagebreak} est intéressante
lorsque utilisée avec son argument optionnel \meta{n}: celui-ci
indique, par le biais d'un nombre entier entre 0 et 4, à quel point
nous \emph{recommandons} à {\LaTeX} d'insérer un saut de page à
l'endroit où la commande apparaît (0 étant une faible recommandation
et 4, une forte).

La commande \cmd{\enlargethispage}, comme son nom l'indique, permet
d'allonger une page de \meta{longueur} pour y faire tenir plus de
texte. C'est une commande particulièrement utile pour éviter que la
toute dernière ligne d'un chapitre ou d'un document se retrouve seule sur
une page.

\subsection{Coupure de mots}
\label{sec:trucs:controle:coupure}

La coupure automatique des mots en fin de ligne est toujours active
avec {\LaTeX}. C'est d'ailleurs pourquoi il importe d'indiquer à
{\LaTeX} dans quelle langue est le texte (lorsque ce n'est pas en
anglais) avec les commandes du paquetage \pkg{babel}.

Il existe deux façons de contrôler la coupure de mots. La première,
principalement utilisée lorsque {\LaTeX} refuse de couper un mot en
fin de ligne, consiste à insérer des \emph{suggestions} d'endroits où
couper le mot avec la commande \cmd{\-}. Par exemple, en écrivant
\verb=vrai\-sem\-blance=, on indique à {\LaTeX} qu'il est possible de
diviser le mot en \emph{vrai-semblance} ou \emph{vraisem-blance}.

La seconde méthode, celle-là surtout utilisée lorsque {\LaTeX} ne
reconnaît pas des mots qui reviennent souvent dans le document,
consiste à fournir dans le préambule une liste d'exceptions avec la
commande
\begin{lstlisting}
\hyphenation`\marg{liste}'
\end{lstlisting}
La \meta{liste} est une suite de mots, séparés par des virgules, des
blancs ou des retours à la ligne, dans lesquels les points de coupure
sont identifiés par un trait d'union.

\begin{exemple}
  La commande suivante, insérée dans le préambule, permet d'ajouter
  des points de coupure aux mots «puisque», «constante» et
  «vraisemblance» pour l'ensemble du document.
\begin{lstlisting}
\hyphenation{puis-que,cons-tante,vrai-sem-blance}
\end{lstlisting}
  \qed
\end{exemple}

\begin{important}
  Règle générale, garder les opérations d'ajustements de la mise en
  page --- position des éléments flottants, sauts de page, lignes trop
  longues, etc. --- pour la toute fin de la rédaction.
\end{important}


\section{Au-delà de la police Computer Modern}
\label{sec:trucs:polices}

{\CM%
  Les documents {\LaTeX} standards sont facilement reconnaissables par
  leur police par défaut, Computer Modern --- celle utilisée dans ce
  paragraphe. Pour qui souhaitait briser la relative monotonie induite
  par cette uniformité, il a longtemps été difficile d'utiliser une
  autre police. Fort heureusement, la situation a beaucoup évolué et
  il est aujourd'hui assez simple de produire des documents {\LaTeX}
  utilisant des polices variées.}

Avant d'aller plus loin, une mise en garde: si un document contient
plus que quelques équations mathématiques très simples, le choix de
police devient très restreint. La raison: peu de polices comprennent
des symboles mathématiques et les informations nécessaires pour les
assembler selon les hauts standards qualité usuels de {\LaTeX}.

Cela dit, pour qui souhaite aller au-delà de la police Computer Modern
sans trop se compliquer la vie, il existe deux solutions principales.

\begin{enumerate}
\item Utiliser l'une ou l'autre des polices PostScript standards
  convenant pour du texte normal
  ({\Times Times}, %
   {\Palatino Palatino}, %
   {\Bookman Bookman}, %
   {\Charter Charter}, %
   {\NewCent New Century Schoolbook}, %
   {\Helvet Helvetica}). %
  C'est très simple avec toute distribution {\LaTeX} moderne: il
  suffit de charger le paquetage approprié. Consulter la %
  \doc{psnfss2e}{http://texdoc.net/pkg/psnfss/} %
  de l'ensemble de paquetages PSNFSS.
\item Utiliser une police OpenType ou TrueType présente sur son
  système avec le moteur {\XeLaTeX}. Seule une poignée de ces polices
  offre toutefois un support approprié pour les mathématiques. La
  gestion des polices avec {\XeLaTeX} se fait avec le paquetage
  standard \pkg{fontspec}; consulter sa %
  \doc{fontspec}{http://texdoc.net/pkg/fontspec/}.
\end{enumerate}

Pour les thèses et mémoires de l'Université Laval, la Faculté des
études supérieures et postdoctorales accepte les polices %
\begin{quote}
  \begin{tabbing}
    Computer Modern \qquad \=  ABCDEF abcdef 1234567890 \kill
    {\CM Computer Modern} \> {\CM ABCDEF abcdef 1234567890} \\
    {\Times Times} \> {\Times ABCDEF abcdef 1234567890} \\
    {\Palatino Palatino} \> {\Palatino ABCDEF abcdef 1234567890}
  \end{tabbing}
\end{quote}
Pour utiliser ces deux dernières avec pdf{\LaTeX}, charger
respectivement les paquetages \pkg{mathptmx} ou \pkg{mathpazo}. Avec
{\XeLaTeX}, utiliser les polices Termes et Pagella du projet %
\link{http://www.gust.org.pl/projects/e-foundry/tex-gyre/}{TeX~Gyre}.
Ce sont des polices très similaires à Times et Palatino, disponibles
en version OpenType et qui fournissent un bon support pour les
mathématiques via le projet frère %
\link{http://www.gust.org.pl/projects/e-foundry/tg-math/}{TeX~Gyre
  Math}.

\begin{exemple}
  \label{ex:trucs:palatino}
  Pour utiliser la police PostScript classique Palatino avec {\LaTeX}
  tant pour le texte que pour les mathématiques, il suffit d'insérer
  dans le préambule de son document la commande
\begin{lstlisting}
\usepackage{mathpazo}
\end{lstlisting}

  Avec le moteur {\XeLaTeX}, il est possible d'utiliser n'importe
  quelle police OpenType et TrueType installée dans le système
  d'exploitation de l'ordinateur. Pour obtenir un résultat équivalent
  à celui de \pkg{mathpazo}, installer les polices TeX~Gyre dans le
  système, puis insérer dans le préambule les commandes
\begin{lstlisting}
\usepackage{fontspec}
\setmainfont{TeX Gyre Pagella}
\setmathfont{TeX Gyre Pagella Math}
\end{lstlisting}
  \qed
\end{exemple}

Le texte principal du présent document est en %
\link{http://tug.org/store/lucida/}{Lucida Bright~OT}, %
une police commerciale de très haute qualité offrant également un
excellent support pour les mathématiques. Ses auteurs ont toujours été
proches de la communauté {\LaTeX}. La Bibliothèque de l'Université
Laval détient une licence d'utilisation de cette police. Les étudiants
et le personnel de l'Université peuvent s'en procurer une copie
gratuitement en écrivant à
\href{mailto:lucida@bibl.ulaval.ca}{lucida@bibl.ulaval.ca}.



\section{Couleurs}
\label{sec:trucs:couleurs}

L'utilisation de couleur dans un document {\LaTeX} requiert de charger
le paquetage \pkg{xcolor} \citep{xcolor}. Celui-ci définit d'abord
plusieurs couleurs que l'on peut utiliser directement; consulter la %
\doc{xcolor}{http://texdoc.net/pkg/xcolor} %
pour en connaître les différentes listes. Le
\autoref{tab:trucs:couleurs} fournit celle des couleurs toujours
disponibles.

\begin{table}
  \centering
  \caption{Couleurs toujours disponibles lorsque le paquetage
    \pkg{xcolor} est chargé}
  \label{tab:trucs:couleurs}
  \begin{tabularx}{1.0\linewidth}{XlXll}
    \toprule
    \fcolorbox{black}{black}{\phantom{xx}}\, black &
    \fcolorbox{black}{blue}{\phantom{xx}}\, blue &
    \fcolorbox{black}{brown}{\phantom{xx}}\, brown &
    \fcolorbox{black}{cyan}{\phantom{xx}}\, cyan &
    \fcolorbox{black}{darkgray}{\phantom{xx}}\, darkgray \\
    \addlinespace[3pt]
    \fcolorbox{black}{gray}{\phantom{xx}}\, gray &
    \fcolorbox{black}{green}{\phantom{xx}}\, green &
    \fcolorbox{black}{lightgray}{\phantom{xx}}\, lightgray &
    \fcolorbox{black}{lime}{\phantom{xx}}\, lime &
    \fcolorbox{black}{magenta}{\phantom{xx}}\, magenta \\
    \addlinespace[3pt]
    \fcolorbox{black}{olive}{\phantom{xx}}\, olive &
    \fcolorbox{black}{orange}{\phantom{xx}}\, orange &
    \fcolorbox{black}{pink}{\phantom{xx}}\, pink &
    \fcolorbox{black}{purple}{\phantom{xx}}\, purple &
    \fcolorbox{black}{red}{\phantom{xx}}\, red \\
    \addlinespace[3pt]
    \fcolorbox{black}{teal}{\phantom{xx}}\, teal &
    \fcolorbox{black}{violet}{\phantom{xx}}\, violet &
    \fcolorbox{black}{white}{\phantom{xx}}\, white &
    \fcolorbox{black}{yellow}{\phantom{xx}}\, yellow \\
    \bottomrule
  \end{tabularx}
\end{table}

Un peu comme pour les changements d'attributs de police, il existe
deux commandes pour modifier la couleur du texte:
\begin{lstlisting}
\color`\marg{nom}'
\textcolor`\marg{nom}\marg{texte}'
\end{lstlisting}
La première modifie la couleur de tout ce qui suit (à moins d'en
limiter la portée avec des accolades) et la seconde, seulement pour
\meta{texte}. Dans les deux cas, \meta{nom} est le nom d'une couleur.
\begin{demo}
  \begin{texample}
\begin{lstlisting}
texte {\color{red} en rouge}
et \textcolor{blue}{en bleu}
\end{lstlisting}
  \producing
  texte {\color{red} en rouge} et \textcolor{blue}{en bleu}
  \end{texample}
\end{demo}

La commande \cmd{\definecolor} permet de définir de nouvelles couleurs
selon plusieurs systèmes de codage. Le plus usuel demeure \emph{Rouge,
  vert, bleu} (RVB ou RGB, en anglais) où une couleur est représentée
par une combinaison de teintes --- exprimées par un nombre entre $0$
et $1$ --- de rouge, de vert et de bleu. Dans ce cas, la syntaxe de
\cmd{\definecolor} est
\begin{lstlisting}
\definecolor`\marg{nom}'{rgb}`\marg{valeur\_r,valeur\_v,valeur\_b}
\end{lstlisting}
où \meta{valeur\_r}, \meta{valeur\_v} et \meta{valeur\_b} sont
respectivement les teintes de rouge, de vert et de bleu.

\begin{exemple}
  La commande
\begin{lstlisting}
\definecolor{acier}{rgb}{0.3,0.4,0.6}
\end{lstlisting}
  définit une nouvelle couleur nommée \code{acier} composée de rouge
  30~\%, de vert 40~\% et de bleu 60~\%: %
  \fcolorbox{black}[rgb]{0.3,0,0}{\phantom{xx}} $+$ %
  \fcolorbox{black}[rgb]{0,0.4,0}{\phantom{xx}} $+$ %
  \fcolorbox{black}[rgb]{0,0,0.6}{\phantom{xx}} $=$ %
  \fcolorbox{black}[rgb]{0.3,0.4,0.6}{\phantom{xx}}. %
  Une fois la couleur  \code{acier} définie dans le préambule, elle
  devient disponible pour être utilisée dans les commandes
  \cmdprint{\color} et \cmdprint{\textcolor}. %
  \qed
\end{exemple}

La commande \cmd{\colorlet}, dont la syntaxe simplifiée est
\begin{lstlisting}
\colorlet`\marg{nom}\marg{couleur}'
\end{lstlisting}
permet de faire référence à la \meta{couleur} déjà existante par
\meta{nom}. C'est pratique pour assigner un nom sémantique à une
couleur.



\section{Hyperliens et métadonnées de documents PDF}
\label{sec:trucs:hyperliens}

Nous en avons déjà traité à quelques reprises, notamment à la
\autoref{sec:organisation:renvois}: le paquetage \pkg{hyperref}
\citep{hyperref} permet de transformer les références dans le
texte en hyperliens cliquables lorsque le document est produit avec
pdf{\LaTeX} ou {\XeLaTeX}. C'est très pratique lors de la consultation
électronique d'un document.

Le paquetage offre une multitudes d'options de configuration; nous
n'en présenterons que quelques unes. On accède aux options de
configuration de \pkg{hyperref} via la commande \cmd{\hypersetup} dans
le préambule. Celle-ci prend en arguments des paires
\code{option=valeur} séparées par des virgules.

Une des principales choses que l'on pourra souhaiter configurer, c'est
le comportement et la couleur des divers types d'hyperliens. On
trouvera ci-dessous les options de configuration pertinentes, leur
valeur (avec en gras la valeur par défaut) ainsi qu'une brève
explication de chacune.

\begin{table}[h]
  \begin{tabularx}{1.0\linewidth}{@{}p{7em}p{6em}X@{}}
    \code{colorlinks} & \code{true}|\code{\textbf{false}} & colorer les
                                                            liens selon
                                                            leur type \\
    \code{linktocpage} & \code{true}|\code{\textbf{false}} & faire du
                                                             numéro de
                                                             page
                                                             plutôt que
                                                             du titre l'hyperlien dans
                                                             la table
                                                             des
                                                             matières \\
    \code{linkcolor} & \meta{couleur} & couleur des liens internes \\
    \code{urlcolor}  & \meta{couleur} & couleur des URL externes \\
    \code{citecolor} & \meta{couleur} & couleur des citations \\
    \code{allcolor}  & \meta{couleur} & couleur pour tous les types d'hyperliens
  \end{tabularx}
\end{table}

Lorsque la valeur admissible d'une option est \code{true} ou
\code{false}, sa seule mention équivaut à \code{true}. La valeur
\meta{couleur} est le nom d'une couleur telle que définie par
\pkg{xcolor} (\autoref{sec:trucs:couleurs}).

Les fichiers PDF peuvent contenir diverses métadonnées sur leur
contenu. Le paquetage \pkg{hyperref} permet d'en définir certaines,
notamment le titre ou l'auteur du document.

\begin{table}[!h]
  \begin{tabularx}{1.0\linewidth}{@{}p{7em}p{6em}X@{}}
    \code{pdftitle}  & texte & titre du document PDF \\
    \code{pdfauthor} & texte & auteur du document PDF
  \end{tabularx}
\end{table}

Consulter la section~3.7 de la %
\doc{hyperref}{http://texdoc.net/pkg/hyperref} %
de \pkg{hyperref} pour obtenir la liste complète des options de
configuration des métadonnées.

\begin{exemple}
  \label{ex:trucs:couleurs}
  Le présent document fait appel aux définitions de couleurs
  et aux options de configuration de \pkg{hyperref} suivantes:
\begin{lstlisting}
\definecolor{link}{rgb}{0,0.4,0.6}   % ~RoyalBlue de dvips
\definecolor{url}{rgb}{0.6,0,0}      % rouge foncé
\definecolor{citation}{rgb}{0,0.5,0} % vert foncé
\hypersetup{%
  pdfauthor = {Vincent Goulet},
  pdftitle = {Rédaction avec LaTeX},
  colorlinks = true,
  linktocpage = true,
  urlcolor = url,
  linkcolor = link,
  citecolor = citation}
\end{lstlisting}
  \qed
\end{exemple}

\begin{exemple}
  Les gabarits de thèses et de mémoires livrés avec la classe
  \class{ulthese} contiennent dans le préambule la ligne suivante:
\begin{lstlisting}
\hypersetup{colorlinks,allcolors=ULlinkcolor}
\end{lstlisting}
  Tous les liens de la thèse ou du mémoire seront donc de la même
  couleur, \code{ULlinkcolor} \fcolorbox{black}{ULlinkcolor}{\phantom{xx}},
  une couleur définie par la classe. %
  \qed
\end{exemple}

\begin{conseil}
  L'interaction du paquetage \pkg{hyperref} avec les autres est
  parfois (voire souvent) délicate. Pour cette raison, il est
  fortement recommandé que \pkg{hyperref} soit le tout dernier
  paquetage chargé dans le préambule.
\end{conseil}



\section{Présentation de code informatique}
\label{sec:trucs:listings}

L'environnement standard \Ie{verbatim} de {\LaTeX} permet de présenter
du texte tel qu'il est entré dans le code source du document. C'est un
environnement particulièrement utile pour afficher du code
informatique puisque le texte est composé en police non
proportionnelle et que sa disposition exacte est respectée.

\begin{demo}
  \begin{texample}
\begin{lstlisting}[deletetexcs={int,include}]
\begin{verbatim}
/* Hello World en C */
#include <stdio.h>

int main()
{
    printf("Hello world\n");
    return 0;
}
\end{verbatim}
\end{lstlisting}
    \producing
\begin{verbatim}
/* Hello World en C */
#include <stdio.h>

int main()
{
    printf("Hello world\n");
    return 0;
}
\end{verbatim}
  \end{texample}
\end{demo}

Si un document doit contenir beaucoup de code informatique et que l'on
souhaite exercer un fin contrôle sur sa disposition et sa mise en
forme, il vaut mieux se tourner vers un paquetage spécialisé comme
\pkg{listings} \citep{listings}. La %
\doc{listings}{http://texdoc.net/pkg/listings} %
du paquetage compare ses fonctionnalités à celles de plusieurs autres
paquetages similaires.

Le paquetage \pkg{listings} peut effectuer automatiquement le marquage
des mots-clés de nombreux langages de programmation, ajouter des
numéros de ligne, importer du code de fichiers externes ou même
indexer les mots-clés des extraits de code. À titre d'illustration, en
utilisant l'environnement \Ie{lstlisting} de \pkg{listings} plutôt que
\Ie{verbatim}, l'extrait de code C ci-dessus pourrait se présenter
ainsi:
\begin{demo}
  \begin{texample}
    \begin{vglisting}
\begin{lstlisting}
/* Hello World en C */
#include <stdio.h>

int main()
{
  printf("Hello world\n");
  return 0;
}
\end{lstlisting}
    \end{vglisting}
    \producing
\begin{lstlisting}[language=C,%
  deletetexcs={int,include},
  frame=single,%
  backgroundcolor=\color{Azure2},%
  rulecolor=\color{black},%
  numbers=left,%
  numberstyle=\tiny\sffamily,%
  stringstyle=\color{orange}\itshape,%
  identifierstyle=\color{cyan}\mdseries,%
  xleftmargin=12pt,%
  keywordstyle=\color{blue}\bfseries,
  index=]
/* Hello World en C */
#include <stdio.h>

int main()
{
  printf("Hello world\n");
  return 0;
}
\end{lstlisting}
  \end{texample}
\end{demo}

Il serait trop long et nettement hors de la portée du présent ouvrage
d'expliquer les nombreuses fonctionnalités de \pkg{listings}.
Précisons simplement que le paquetage a servi pour en composer les
extraits de code et pour construire une grande partie de
l'index.

\begin{exemple}
  \label{ex:trucs:listings}
  Pour parvenir à la présentation des extraits de code source {\LaTeX}
  de ce document, le paquetage \pkg{listings} est configuré dans le
  préambule de la manière suivante:
\begin{lstlisting}
%% Couleurs
\definecolor{comments}{rgb}{0.7,0,0}

%% Configuration de listings
\lstset{language = [LaTeX]TeX,
  basicstyle = \ttfamily\NoAutoSpacing,
  keywordstyle = \mdseries,
  commentstyle = \color{comments}\slshape,
  extendedchars = true,
  showstringspaces = false,
  backgroundcolor = \color{LightYellow1},
  frame = lr,
  rulecolor = \color{LightYellow1},
  xleftmargin = 3.4pt,
  xrightmargin = 3.4pt}
\end{lstlisting}
  (La couleur \code{LightYellow1} est définie par \pkg{xcolor} lorsque
  le paquetage est chargé avec l'option \code{x11names}.)
  \qed
\end{exemple}



\section{Production de rapports avec l'analyse intégrée}
\label{sec:trucs:sweave}

Les publications scientifiques reposent souvent sur une forme ou une
autre d'analyse numérique ou statistique, la production de code
informatique, une simulation stochastique, etc. La portion
développement et analyse est alors produite avec un certain outil
et la publication, avec un outil d'édition séparé --- {\LaTeX} dans le cas
qui nous occupe. Or, tous les auteurs ont vécu cette situation: les
résultats de l'analyse changent et il faut modifier le rapport en
conséquence, refaire les tableaux et les graphiques, retracer cette
valeur isolée au fil du texte directement tirée de l'analyse\dots\
Seule la quantité de temps perdu rivalise avec le risque d'erreur.

Il existe pourtant une meilleure façon de travailler.

Cette meilleure façon de faire, tirée du concept de
\emph{programmation lettrée}, consiste à combiner dans un seul et même
document l'analyse et le rapport, puis de produire automatiquement une
partie du second à partir de la première.

Les utilisateurs du système statistique R bénéficient d'une mise en
œuvre simple et élégante du concept ci-dessus avec l'outil Sweave
\citep{Sweave}. Un fichier Sweave est à la base un document {\LaTeX}
dans lequel on a inséré du code R à l'intérieur de balises spéciales
\verb|<<>>=| et \verb|@| tirée de la syntaxe \code{noweb}
\citep{noweb}, comme ceci:
\begin{lstlisting}
\section{Commandes R}

L'utilisateur de R interagit avec l'interprète en entrant
des commandes à l'invite de commande:
<<>>=
2 + 3
pi
cos(pi/4)
@
La commande \verb=exp(1)= donne \Sexpr{exp(1)},
la valeur du nombre $e$.
\end{lstlisting}

Par convention, on enregistre un tel document sous un nom se terminant
par l'extension \code{.Rnw}. Sa compilation s'effectue
en deux étapes:
\begin{enumerate}
\item Le fichier \code{.Rnw} est passé à la commande \code{Sweave()}
  de R. Celle-ci retrace les extraits de code et les remplace par des
  environnements {\LaTeX} contenant, par défaut, les expressions R et
  leur résultat dans des environnements \Ie{Sinput} et \Ie{Soutput}.
  Elle évalue également les expressions R se trouvant dans les
  commandes \cmd{\Sexpr} pour les remplacer par leur résultat. Cela
  produit un fichier \code{.tex}:
\begin{lstlisting}
\section{Commandes R}

L'utilisateur de R interagit avec l'interprète en entrant
des commandes à l'invite de commande:
\begin{Schunk}
\begin{Sinput}
> 2 + 3
\end{Sinput}
\begin{Soutput}
[1] 5
\end{Soutput}
\begin{Sinput}
> pi
\end{Sinput}
\begin{Soutput}
[1] 3.141593
\end{Soutput}
\begin{Sinput}
> cos(pi/4)
\end{Sinput}
\begin{Soutput}
[1] 0.7071068
\end{Soutput}
\end{Schunk}
La commande \verb=exp(1)= donne 2.71828182845905,
la valeur du nombre $e$.
\end{lstlisting}
\item On compile le fichier \code{.tex} comme d'habitude.
\end{enumerate}

Sweave se révèle particulièrement utile pour créer des graphiques à
partir de R: tout ce que l'on doit conserver dans son fichier
\code{.Rnw}, c'est le code pour créer le graphique. Il est également
possible de contrôler l'exécution des blocs de code et l'affichage du
code source et des résultats par le biais d'options placées à
l'intérieur de la balise d'ouverture \verb|<<>>=|. Consulter la %
\link{http://leisch.userweb.mwn.de/Sweave/Sweave-manual.pdf}{documentation} %
de Sweave pour les détails.

\link{http://mpastell.com/pweave/}{Pweave} est un système similaire à
Sweave pour le langage Python.

Inspiré de Sweave, knitr \citep{knitr} permet également d'entrelacer
du code {\LaTeX} et du code R. Cet outil offre plus d'options de
traitement que Sweave, mais au prix d'une complexité accrue.

En terminant, rappelons que Sweave n'est qu'un exemple de système de
programmation lettrée. Il en existe plusieurs autres. Il appartient
dès lors au lecteur de trouver le système qui correspond le mieux à
ses besoins.

\begin{information}
  On doit le concept de programmation lettrée au créateur de {\TeX},
  Donald Knuth.  En fait, tout le code source de {\TeX} est écrit
  en programmation lettrée! La %
  \link{https://fr.wikipedia.org/wiki/Programmation_lettrée}{page Wikipedia} %
  consacrée au sujet offre un très bon survol de l'historique et de la
  nature du concept.

  Le lecteur intéressé pourra consulter, par exemple, le %
  \link{http://www.ctan.org/pkg/ulthese}{code source} %
  de la classe \class{ulthese}. La documentation de la classe, le
  code {\LaTeX}, les gabarits, tout se trouve entrelacé
  dans le seul fichier \code{ulthese.dtx}.
\end{information}


\section{Diapositives}
\label{sec:trucs:diapositives}

Il n'est pas rare qu'une publication scientifique fasse l'objet d'une
présentation dans le cadre d'un colloque ou d'un séminaire. Lorsque le
texte a été rédigé avec {\LaTeX}, il est tout naturel de souhaiter le
réutiliser pour la préparation de diapositives --- surtout si le texte
comporte de nombreuses équations mathématiques qu'il serait
extrêmement long de retranscrire dans un logiciel de présentation
comme PowerPoint.

Fort heureusement, il est tout à fait possible de composer ses
diapositives avec {\LaTeX}. La classe standard \class{slides} produit
des diapositives élégantes, quoique minimalistes, qui conviendront à
l'auteur qui ne recherche rien de plus que du texte noir sur fond
blanc.

Cependant, l'outil devenu le standard \emph{de facto} pour la
production de diapositives est la classe \class{beamer}
\citep{beamer}. Celle-ci compte un grand nombre de thèmes et de
gabarits élaborés, rend très simple l'insertion d'animations d'une
diapositive à l'autre, gère automatiquement la table des matières
et\dots\ produit des diapositives en couleur.

Quelle que soit la classe utilisée, les diapositives produites avec
{\LaTeX} se présentent sous forme de fichier PDF. On les projette avec
une liseuse PDF en mode plein écran.

Nous n'irons pas plus loin sur le sujet dans ce document. Consulter
la %
\doc{beamer}{http://texdoc.net/pkg/beameruserguide} %
de \class{beamer} pour apprendre à utiliser la classe. Produire des
diapositives de grande qualité avec les gabarits fournis avec la
classe se révèle simple et rapide.



\section{Gestion des versions et travail collaboratif}
\label{sec:trucs:cvs}

Plusieurs personnes travaillent sur un même fichier, ou encore, une
seule personne y travaille, mais de plusieurs postes de travail
différents. Quelle est la plus récente version du fichier? Un ajout
fait hier dans le fichier a-t-il été pris en compte par une collègue
aujourd'hui? Une modification apportée au fichier n'est plus
nécessaire; comment retourner en arrière facilement?

Les informaticiens ont résolu ce genre de problèmes il y a des
dizaines d'années avec les systèmes de gestion de versions. Les
systèmes les plus populaires en ce moment sont Subversion
\citep{subversion} et Git \citep{git}.

Bien que développés à l'origine pour la gestion du code source de
logiciels, les systèmes de contrôle de versions conviennent
parfaitement pour les sources {\LaTeX}. L'utilisation d'un tel système
permet de:
\begin{itemize}
\item toujours savoir quelle est la plus récente version d'un fichier;
\item travailler à plusieurs personnes simultanément sur un même
  fichier;
\item revenir aisément à une version antérieure d'un fichier;
\item comparer deux versions d'un fichier pour connaître les
  modifications qui y ont été apportées;
\item gérer automatiquement les éventuels conflits de modification
  d'un fichier;
\item disposer en tout temps d'une copie de secours de son travail
  lorsque l'on a recours à un serveur central (obligatoire avec
  Subversion; optionnel avec Git\footnote{%
    Si ce n'est que pour le volet de sauvegarde, nous recommandons
    d'utiliser avec Git un dépôt central tel que
    \link{https://github.com}{GitHub}, un serveur public qui héberge
    déjà des millions de projets.}.)
\end{itemize}
Un système de gestion de versions est un outil qui permet d'augmenter
considérablement sa productivité ou celle de son équipe de
travail au moment de rédiger un ouvrage scientifique.

La %
\link{https://fr.wikipedia.org/wiki/Gestion_de_versions}{page Wikipedia} %
sur le sujet offre une bonne introduction aux systèmes de gestion de
versions. Les membres de la communauté de l'Université Laval qui
souhaitent mettre sur pied un dépôt pour un projet pourront
consulter leur équipe de soutien informatique facultaire.

\begin{information}
  La présente documentation est conservée dans un dépôt Subversion
  public d'où l'on peut toujours obtenir la plus récente version.
  Consulter la page des notices de copyright au début du document pour
  accéder au dépôt.
\end{information}



%%%
%%% Exercices
%%%

\section{Exercices}
\label{sec:trucs:exercices}

\Opensolutionfile{solutions}[solutions-trucs+astuces]

\begin{Filesave}{solutions}
\section*{Chapitre \ref*{chap:trucs}}
\addcontentsline{toc}{section}{Chapitre \protect\ref*{chap:trucs}}

\end{Filesave}

\noindent%
Pour les exercices
\nolink{\ref{exercice:trucs:1}}--\nolink{\ref{exercice:trucs:n}},
utiliser le fichier \fichier{exercice\_trucs.tex}. Celui-ci reprend
une partie de la documentation de la classe \class{ulthese}. De plus,
on remarquera que le document:
\begin{itemize}
\item doit être compilé avec pdf{\LaTeX} puisqu'il charge les
  paquetages \pkg{inputenc} et \pkg{fontenc};
\item définit des nouvelles commandes \cmdprint{\class} et
  \cmdprint{\fichier} pour composer, respectivement, les noms de
  classe et les noms de fichier;
\item utilise la commande \cmdprint{\doc} de
  l'\autoref{ex:commandes:doc};
\item charge le paquetage \pkg{hyperref}, ce qui transforme les titres
  de la table des matières, les renvois aux notes de bas de page et
  les liens externes en hyperliens.
\end{itemize}
\medskip

\begin{exercice}
  \label{exercice:trucs:1}
  Compiler et visualiser le fichier sans aucune modification. Le texte
  est composé dans la police par défaut Computer Modern.

  Ensuite, modifier le préambule du document pour composer le document
  dans la police Palatino, tel qu'expliqué à
  l'\autoref{ex:trucs:palatino}. Charger également le paquetage
  \pkg{helvet} afin d'utiliser Helvetica pour le texte en police sans
  empattements (\cmdprint{\textsf}).
  \begin{sol}
    On doit ajouter dans le préambule les commandes
\begin{lstlisting}
\usepackage{mathpazo}
\usepackage{helvet}
\end{lstlisting}
    La police Helvetica est très grande. Tel que mentionné dans la %
    \doc{psnfss2e}{http://texdoc.net/pkg/psnfss/} %
    de PSNFSS (section~4), il est généralement préférable de charger
    le paquetage \pkg{helvet} avec, par exemple,
\begin{lstlisting}
\usepackage[scaled=0.92]{helvet}
\end{lstlisting}
    pour que le texte en Helvetica se marie mieux à celui dans une
    autre police.
  \end{sol}
\end{exercice}

\begin{exercice}
  Configurer le paquetage \pkg{hyperref} pour que les hyperliens dans
  la table des matières soient ancrés aux numéros de page plutôt
  qu'aux titres de section.
  \begin{sol}
    Tel qu'expliqué à la \autoref{sec:trucs:hyperliens}, on doit
    insérer dans le préambule du document la commande
\begin{lstlisting}
\hypersetup{linktocpage=true}
\end{lstlisting}
    ou, plus simplement,
\begin{lstlisting}
\hypersetup{linktocpage}
\end{lstlisting}
  \end{sol}
\end{exercice}

\begin{exercice}
  Charger le paquetage \pkg{xcolor} et ajouter l'option
  \code{colorlinks} à \pkg{hyperref}, puis recompiler le document.
  Observer les changements.
  \begin{sol}
    En conservant l'ajout de l'exercice précédent, on a dans le
    préambule la commande
\begin{lstlisting}
\hypersetup{colorlinks, linktocpage}
\end{lstlisting}
    Les hyperliens se présentent maintenant en couleur selon les
    paramètres par défaut de \pkg{hyperref}.
  \end{sol}
\end{exercice}

\begin{exercice}
  En s'inspirant de l'\autoref{ex:trucs:couleurs}, modifier la couleur
  des liens internes et externes.
  \begin{sol}
    On peut soit utiliser des couleurs prédéfinies de \pkg{xcolor}
    (\autoref{tab:trucs:couleurs}), soit en  définir de nouvelles avec
    \cmd{\definecolor}. On fait ensuite appel aux couleurs choisies
    pour les options \code{linkcolor} (liens internes) et
    \code{urlcolor} (liens externes) de \pkg{hyperref}.

    Exemple utilisant des couleurs prédéfinies:
\begin{lstlisting}
\hypersetup{colorlinks, linktocpage,
  linkcolor=brown, urlcolor=blue}
\end{lstlisting}

    Exemple avec de nouvelles couleurs:
\begin{lstlisting}
\definecolor{link}{rgb}{0,0.4,0.6}
\definecolor{url}{rgb}{0.6,0,0}
\hypersetup{colorlinks, linktocpage,
  urlcolor=url, linkcolor=link}
\end{lstlisting}
  \end{sol}
\end{exercice}

\begin{exercice}
  \label{exercice:trucs:n}
  Charger le paquetage \pkg{listings} et modifier l'environnement
  \Pe{verbatim} que l'on trouve dans le document pour un environnement
  \Pe{lstlisting}. En s'inspirant de l'\autoref{ex:trucs:listings},
  configurer la présentation des extraits de code pour utiliser une
  police non proportionnelle (\cmdprint{\ttfamily}) et un arrière-plan
  de la couleur standard \code{lightgray}.
  \begin{sol}
    S'il y avait plusieurs extraits de code dans le document, mieux
    vaudrait les configurer tous à l'identique dans le préambule du
    document avec
\begin{lstlisting}
\lstset{basicstyle=\ttfamily,
  backgroundcolor=\color{lightgray}}
\end{lstlisting}
    Ensuite,
\begin{vglisting}
\begin{lstlisting}
latex ulthese.ins
\end{lstlisting}
\end{vglisting}
    donne le résultat demandé.

    Pour un seul extrait, il est également possible de simplement
    charger le paquetage dans le préambule et d'effectuer la
    configuration à l'ouverture de l'environnement, comme ceci:
\begin{vglisting}
\begin{lstlisting}[basicstyle=\ttfamily,
  backgroundcolor=\color{lightgray}]
latex ulthese.ins
\end{lstlisting}
\end{vglisting}
  \end{sol}
\end{exercice}

\Closesolutionfile{solutions}


%%% Local Variables:
%%% mode: latex
%%% TeX-engine: xetex
%%% TeX-master: "formation-latex-ul"
%%% coding: utf-8
%%% End:


\appendix
version https://git-lfs.github.com/spec/v1
oid sha256:1ae860479f70bc51be429ddc5cb7a2192ed0d0bd7a682f6d702ab12b30b60fb6
size 16486
               % éléments spécifiques à ulthese
version https://git-lfs.github.com/spec/v1
oid sha256:f0be3950c0a12df00719de0cad03bcc32b9767ee849cb1a718ee232f110f1df8
size 431


\bibliography{formation-latex-ul}

\cleardoublepage
\printindex

\cleartoverso

\thispagestyle{empty}
\vspace*{\fill}

\begingroup
\calccentering{\unitlength}
\begin{adjustwidth*}{\unitlength}{-\unitlength}
  \begin{flushleft}
    \small %
    Ce document a été produit avec le système de mise en page
    {\XeLaTeX} à partir de la classe \class{memoir}. Le texte
    principal est en Lucida Bright~OT 11~points, les mathématiques en
    Lucida Bright Math~OT, le code informatique en Lucida Grande
    Mono~DK et les titres en Adobe Myriad~Pro. Les icônes proviennent
    de la police Font~Awesome.
  \end{flushleft}
\end{adjustwidth*}
\endgroup
\vfill

%%% Local Variables:
%%% mode: latex
%%% TeX-engine: xetex
%%% TeX-master: "formation-latex-ul"
%%% coding: utf-8
%%% End:


\cleartoverso

%% Page couverture arrière.
\pagestyle{empty}
\begingroup

\TPGrid{3}{36}
\textblockorigin{0mm}{0mm}
\setlength{\parindent}{0mm}
\setlength{\banderougewidth}{2\TPHorizModule}
\setlength{\bandeorwidth}{\TPHorizModule}
\setlength{\gapwidth}{2pt}
\addtolength{\bandeorwidth}{-\gapwidth}

%% bandeau identitaire arrière
\begin{textblock*}{8.5in}[0,1](0mm,30\TPVertModule)
  \textcolor{or}{\rule{\bandeorwidth}{\TPVertModule}}%      % bande or
  \rule{\gapwidth}{0pt}%                                    % filet
  \textcolor{rouge}{\rule{\banderougewidth}{\TPVertModule}} % bande rouge
\end{textblock*}

% code-barre
\begin{textblock*}{0.9\TPHorizModule}(0.1\TPHorizModule,25\TPVertModule)
  \includegraphics[height=4\TPVertModule]{codebarre_\ISBN}
\end{textblock*}

\endgroup

%%% Local Variables:
%%% mode: latex
%%% TeX-master: "formation-latex-ul"
%%% TeX-engine: xetex
%%% End:


\end{document}

%%% Local Variables:
%%% mode: latex
%%% TeX-engine: xetex
%%% TeX-master: t
%%% coding: utf-8
%%% End:
